% ============================================================
%  R. Nielsen, "Mr. S. Kierkegaard's »Johannes Climacus« and
%  Dr. H. Martensen's »Christian Dogmatics«. An Examining Review."
%  Copenhagen: Published by University Bookseller C. A. Reitzel;
%  Printed by Royal Court Printer Bianco Luno, 1849.
%  ENGLISH TRANSLATION — from transcription.tex (Danish, complete).
%
%  Method: ../../../TRANSLATION-PLAYBOOK.md. Fill each "% [text to be added:
%  pp. X--Y]" marker via the Edit tool, in order, then run the sandbox
%  compile recipe from the playbook (lmodern/babel substitution).
%
%  STRUCTURE: mirrors transcription.tex 1:1 — a continuous review-essay, no
%  chapters or numbered sections. Two levels of heading:
%    1. Two display questions (large bold), the book's two halves' hinge:
%       p.9  "Has not 'Johannes Climacus' dialectically sharpened the
%             Christian problem?"
%       p.48 "Has not 'the Christian Dogmatics' undialectically evaded the
%             Christian problem?"
%    2. Lettered sub-heads (letterspaced -> \emph{}, centred), TWO series,
%       restarting at the hinge (p.42, centred rule):
%       Climacus half:  a) p.11   b) p.21   c) p.27   d) p.34
%         (inside a): the further split alpha) p.12 / beta) p.17, run-in,
%         typed as $\alpha$)/$\beta$) — no textalpha dependency)
%       Martensen half: a) p.57   b) p.67   c) p.90   d) p.125
%  Batch boundaries below follow the transcription's own 16 batches (see
%  RESUME-NOTES.md), each ~8 printed pages: 3-10, 11-18, 19-29, 30-37,
%  38-46, 47-54, 55-62, 63-70, 71-78, 79-86, 87-94, 95-102, 103-110,
%  111-118, 119-126, 127-132.
%
%  Quote-balance ledger note (from transcription.tex's RESUME-NOTES): the
%  Danish print itself opens quotations it never closes and vice versa, at
%  six documented sites (pp.12/14/18, p.46 fn., pp.59-60, p.71, p.92,
%  pp.103/105, p.123 fn.) — genuine printer's defects, carried into this
%  translation as inline `%` comments at each site, same convention used
%  throughout this corpus.
% ============================================================
\documentclass[12pt]{book}
\usepackage[utf8]{inputenc}
\usepackage[T1]{fontenc}
\usepackage[english]{babel}
\usepackage[protrusion=true,expansion=true]{microtype}
\usepackage[margin=1.4in]{geometry}
\usepackage{setspace}
\usepackage{amsmath}
\usepackage{enumitem}
\usepackage{graphicx}    % for \reflectbox (the old Danish »ɔ:« id-est mark, if quoted)
\DeclareUnicodeCharacter{0254}{\reflectbox{c}}  % ɔ (U+0254) = reversed c
\usepackage{libertinus}
\usepackage{libertinust1math}
\usepackage{textalpha}   % polytonic Greek (pp.63-64, 74-75, 109), copied verbatim
\usepackage{fancyhdr}
\usepackage[colorlinks=true,linkcolor=black,urlcolor=black,
  pdftitle={Mr. S. Kierkegaard's Johannes Climacus and Dr. H. Martensen's Christian Dogmatics},
  pdfauthor={Rasmus Nielsen}]{hyperref}
\setlength{\headheight}{15pt}
\pagestyle{fancy}
\fancyhf{}
\fancyhead[LE,RO]{\thepage}
\fancyhead[RE]{\textit{An Examining Review}}
\fancyhead[LO]{\textit{Kierkegaard's »Johannes Climacus«}}
\setlength{\parindent}{1.5em}
\setlength{\parskip}{0pt}
\setstretch{1.3}

\begin{document}

% ============================================================
% TITLE PAGE
% ============================================================
\begin{titlepage}
\centering
\vspace*{2.0cm}
{\Large Mr. S. Kierkegaard's\par}
\vspace{0.5cm}
{\Huge\bfseries ``Johannes Climacus''\par}
\vspace{1.0cm}
{\Large and\par}
\vspace{1.0cm}
{\Large Dr. H. Martensen's\par}
\vspace{0.5cm}
{\Huge\bfseries ``Christian Dogmatics.''\par}
\vspace{1.6cm}
\rule{4cm}{0.4pt}\par
\vspace{1.6cm}
{\LARGE\bfseries An Examining Review\par}
\vspace{0.6cm}
{\large by\par}
\vspace{0.6cm}
{\LARGE\bfseries R. Nielsen,\par}
\vspace{0.25cm}
{\small Professor of Philosophy.\par}
\vfill
{\large\bfseries Copenhagen.\par}
{\small Published by University Bookseller C. A. Reitzel.\par}
{\footnotesize Printed by Royal Court Printer Bianco Luno.\par}
{\large\bfseries 1849.\par}
\vspace{0.4cm}
{\small (translated from the Danish)\par}
\end{titlepage}

\frontmatter
\mainmatter
\setcounter{page}{3}

% ============================================================
% TEXT — printed pp. 3--132
% ============================================================

This review must begin with an apology, that it takes the liberty of setting side by side two writings so dissimilar in direction, spirit, and tone as:

\begin{center}
{\large\bfseries Concluding Unscientific Postscript\par}
{\large to the Philosophical Fragments.\par}
\vspace{0.35em}
{A Mimic-Pathetic-Dialectic Compilation,\par}
{An Existential Contribution,\par}
{by Johannes Climacus.\par}
\vspace{0.35em}
{\small Published by\par}
{S. Kierkegaard.\par}
\vspace{0.2em}
{\small Copenhagen 1846, C. A. Reitzel.\par}
\vspace{0.9em}
{and\par}
\vspace{0.9em}
{\large\bfseries Christian Dogmatics,\par}
{set forth by\par}
{\bfseries Dr. H. Martensen.\par}
\vspace{0.2em}
{\small Copenhagen 1849, C. A. Reitzel's Press.\par}
\end{center}

% --- p. 4 ---
To set an unscientific postscript alongside a systematic-scientific principal work, an experimenting humorist alongside a Christian dogmatician, a fleeting pseudonym alongside a renowned author, will perhaps easily awaken, in many, a disquieting association of ideas, and produce the same uncomfortable impression as if one were now, in reckless uncriticalness, to confound the serious with the jesting, the ecclesiastical with the theatrical --- a preacher with a solo dancer. But, since this impropriety presses itself upon us so strongly, the reviewer is in earnest concerned, so far as it lies with him, to remove it, and shall now at once give the reader an account of the respect in which the two writings may, indeed, be supposed to throw light on one another, and of the manner in which the dissimilarity with which they so abruptly repel one another can here be --- not, indeed, abolished, but nonetheless, in the form of the juxtaposition, considerably softened.

The question, around which theology and philosophy at the present time must be said to move, as around the true and proper question of principle, can briefly be set forth in the words:

\begin{center}
\emph{Will the truth of Christianity, by its very nature, be an object of objective knowledge?}
\end{center}

That this question, when it is grasped and understood as it ought to be, is of thoroughgoing significance, is easy to see. As a question of principle it can, in no way, be left hanging in evasive indeterminacy,
% --- p. 5 ---
but must be answered unconditionally, with yes or no. If the question is answered with an unconditional yes, then school-theology is well within its rights, when it continues, as hitherto, to declare itself indispensable to the Church. Theology's victory over worldly knowledge is then an essential moment in faith's victory over unbelief, just as, conversely, theology's scientific defeat is not merely a literary misfortune for the theologians, but a painful loss for the true congregation, an intellectual affliction that must fill the believers with sorrow and grief. It is, then, not only on behalf of culture and learning, but in the holy name of faith and Christianity, that one ought, as hitherto, to hold fast both to ``the critical introduction'' and ``hermeneutics'' and ``apologetics'' and ``the grounding, comprehending dogmatics''; indeed, one ought perhaps even, especially in these times, consistently to have a special church-prayer for ``Christian science,'' a general day of prayer for ``the critical introduction,'' that its ``Easter computation'' might rightly strengthen our faith, and theology hold fast to the Church.

If, on the other hand, the question is answered negatively, not with an evasive turn of phrase, as when it is said, ``that faith is the one thing needful, and that those who feel no need to reflect upon their faith may well let it be,'' but with an open, unconditional no; if it is admitted straightforwardly, that the Christian revelation neither can nor will be an object of our objective knowledge: then the matter takes a remarkable turn; then the link snaps, by which school-theology has hitherto hung itself
% --- p. 6 ---
upon the Church; then the Church, too, in this matter, will know to render unto Caesar the things that are Caesar's, leaving the learned disciplines to their own scientific fate.

Here, then, the matter at hand is not, in the first instance, a scientific distinction between the true and the false theology, but the theological science's principled relation to Christianity; nor is the question here one of the difference between the true and the false speculation, but the inquiry concerns the question, whether speculation has, at all, any religious and Christian validity, or whether religion is not, rather, itself transformed into something quite different, when it is made an object of speculative consideration.

By holding fast to the point of view here set forth, one will, without doubt, admit, that the two reviewed writings belong more closely together than might, at first glance, appear; for while ``the Dogmatics'' proceeds from the presupposition, that the content of Christianity ought to be expounded, grounded, and explained in the forms of objective knowledge, and thus, to that extent, though with some uncertainty, answers our question with yes; the pseudonym's ``Unscientific Postscript'' is, surely, in any case, to be understood as an energetic protest against the present age's objective cognition of Christianity, and thus answers the question with an unconditional, decisive no.
% --- p. 7 ---
A considerable difficulty lies, however, in this, that the pseudonym employs, throughout, the indirect communication in doubly reflected categories, and thereby makes it so precarious a matter to render his own train of thought in a simple report; but, in part, there are found in the book itself not a few passages, in which the discourse approaches being ``instructive'' (though what is said undeniably takes on its own particular cast from the humorous context out of which it springs); in part, the report I here intend to give may quite straightforwardly be explained as my own understanding, as the yield I believe myself to have won by studying the book. If, then, it is in this way forestalled, that my understanding of the book's scientific significance be confused with the author's perhaps somewhat divergent main thought; then I may, indeed, also, without distorting the original's content, avail myself, for my part, of a straightforward form of communication, and thereby possibly even gain several advantages. In the first place, I need not, then, enter into the work's artistic design and peculiar articulation, but can content myself with bringing out the points that most nearly concern the question of principle. In the second place, I can, without departing too far from the author's own language, here and there modify the expression, so that the ``piquant,'' the ``restless,'' the ``unrestrained,'' the ``desultory'' is not permitted to work disturbingly upon the serious reader.\footnote{For the same reason I do not consider it appropriate, either, to indicate paragraph and page for the individual citations, as one otherwise does when citing from a source-work; for, surely, this Postscript is, in its kind, a source-work, but precisely because it is a source-work that is best read \textit{in optima forma}, it may, indeed, best be left to the reader himself to find the cited passages in their original context,
% --- footnote continues on p. 8 ---
whereby he will, at the same time, be able to have the pleasure of laughing and thinking at once.} Finally
% --- p. 8 ---
the whole matter has yet another side, which one might well wish to see somewhat more closely clarified. For it looks, indeed, almost as though theological opinion would now turn in such a way, that the objections which the dialecticizing pseudonym has raised against the objective treatment of ethics and Christianity should deserve no attention whatsoever from ``literature,'' but should exist only for ``a few individuals, who, without feeling any drive toward coherent thinking, can satisfy themselves by thinking in scattered thoughts and aphorisms, whims and flashes.'' In order, then, if possible, to find out whether this opinion is well-grounded, I shall, in all brevity, attempt to set forth the relevant main points in the connection in which they have confirmed themselves to me. Should the attempt fail, should the connection break down, should the underlying thought be weighed and found wanting; then the blame cannot, naturally, fall back upon a ``Postscript,'' which announces itself as ``unscientific,'' but the reproach of criticism falls upon me, who have here believed myself to see something of universal validity in that which, more closely examined, is, after all, only an amusing whim, a whimsical outpouring of a single, arbitrary subjectivity. If, on the other hand, the presentation succeeds even only somewhat, so that it can, at least, be discerned, that a principle is in motion; then the credit is not mine, but the humorous
% --- p. 9 ---
author's, who, after being deprived of his personal freedom and his brilliant weapons, and forced into the stiff form of an instructive discourse, can still hold out in a contest with objective knowledge. Let the outcome be as it will, I conclude as follows: if the pseudonym has no principle, then his wit is ill applied, and ``science'' had then best tell him plainly, that he must not, in the end, go about and addle our studying youth. But if there is something in what he says, if he has really set a question of principle in motion: then, too, one cannot well let him pass with the remark, that he is merely an individuality, a gifted eccentric, going about as an exception; for a principle is more than an individual, and ``science'' dare not let a principle be lost among the exceptions.

On this risk, then, must the review take its course, and the reader take its various moments as contributions toward an answer to the question:

\begin{center}
{\large\bfseries Has not ``Johannes Climacus'' dialectically sharpened
the Christian problem?\par}
\end{center}

If there is to be talk of a Christian-ecclesiastical science: then it is, indeed, evident, that this, like every other science, must aim at a solution of certain definite problems. But in order to solve a problem, one must first know where it is, and what it is, and know how to set it forth with logical sharpness. A scientific presentation, however Christian-ecclesiastically minded it may otherwise be, will, when it betrays that its author
% --- p. 10 ---
has either not known the problem at all, or has, at any rate, seen it only as in a hovering conception, be \textit{eo ipso} a weak piece of work. In theological literature, where all manner of learned inquiries so easily draw attention away from the proper point of view, the true critical ground-point at which faith and knowledge meet, it is of particular importance to be reminded of this. In mathematics, natural science, and history the problems may perhaps show themselves in a more immediate way; but it is certain, that in theology and philosophy one has almost the same labor in stating the problem as in solving it. In these sciences the essential task can, then, at times, become so clouded and obscured, that the thinker is already exhausted and spent, before he has worked his way through to the question of principle, and discovered the turning point from which the whole development should now, in fact, first properly begin. Difficult it has hitherto been; still more difficult it becomes, the more the thetic presuppositions are supplanted by the dialectical. In an immediate, thetic form one can still, more or less, distinguish philosophy from the doctrine of faith; but, once dialectic begins to make the boundaries fluid: when it is said, for instance, that the difference between the supernatural and the natural is fluid, the difference between the subjective and the objective fluid, the difference between explicative and speculative development of concepts fluid, and the difference between the comprehensible and the incomprehensible likewise fluid; then great acuity and much caution is needed,
% --- p. 11 ---
to prevent all this fluidity from finally flowing together into one.

The difficulty is here so much the greater, in that one cannot, however, straightforwardly return to the thetic separation, since dialectic is undeniably the higher form of the movement of ideas. What is required here, then, is not merely a new distinction, but a distinction of a higher kind, a dialectical tension, under which the problem again comes into view: what is required here is that dialectic be corrected by dialectic. With this presupposition, we turn, then, now to the pseudonym, in order to see:

\begin{center}
\emph{a) How the problem arises.}
\end{center}

Without himself having yet entered into Christianity, a reasonable human being can, nonetheless, well understand this much, ``that it will make the single individual eternally blessed, and that it presupposes, precisely in the single individual, this infinite interest in his blessedness, as \textit{conditio sine qua non}.'' The question is, then, not, in the first instance, about the truth of Christianity in and for itself, but about ``the individual's relation to Christianity.''

If, now, the cognition of the truth of Christianity rests exclusively upon the individual's true relation to Christianity, then the failed endeavor will, likewise, be recognizable by the untrue relation, that is, by the misrelation. If it can, then, be proven, that every attempt to secure for oneself the truth of Christianity by the objective way must end in an unavoidable misrelation between the individual and Christianity; then it is, indeed, at the same time proven, that the objective way leads away from
% --- p. 12 ---
the problem. --- \emph{The objective way divides into $\alpha$) the historical and $\beta$) the speculative.}

$\alpha$) If Christianity is regarded as a ``historical document,'' and the inquiring subject is infinitely interested in his relation to Christianity, then the task will, indeed, be to obtain wholly reliable information as to what, then, actually is the Christian doctrine. This task must necessarily bring the subject to despair, since nothing is easier to see, than that the greatest possible certainty in relation to the historical is only an approximation, and that such an approximation is too little to build one's blessedness upon, indeed, ``so dissimilar to an eternal blessedness, that no sum total can result.''

When, thus, the Holy Scripture is regarded as the sure refuge, that decides what is Christian and what is not, then it becomes a matter of securing the Scripture historically-critically. ``One treats here, then: the individual writings' belonging to the canon, their authenticity, integrity, the trustworthiness of their authors, and one posits a dogmatic guarantee: inspiration. When one thinks of the Englishmen's labor on the tunnel, the immense exertion of strength, and how then a small accident can, for a long time, disturb the whole --- then one gets an appropriate notion of this whole critical undertaking. What time, what diligence, what splendid abilities, what excellent learning have here, from generation to generation, been requisitioned for this wonder-work. And yet a little dialectical doubt, touching the presuppositions, can here suddenly
% --- p. 13 ---
disturb the whole, disturb the underground way to Christianity, which one has wished to construct objectively and scientifically, instead of letting the problem come to be as it is, subjectively.''

To secure the Bible by help of the Church will succeed just as little, as, conversely, to secure the Church by help of the Bible, when the matter is kept purely objective. For the present, however, this much is gained, that the Church, with the living Word, the confession of faith, and the sacraments, is a presence, while the Scripture remains in the past. ``The Church,'' it is said, ``cuts away all that proving and proving
% Printer's defect, carried over from transcription.tex: the print
% genuinely repeats the word "proving" here (Danish "Bevisen og
% Bevisen"), zoom-verified there and reproduced verbatim.
that was needed in relation to the Bible. To demand a proof from the Church that it exists is nonsense, just as to demand of a living man proof that he exists. The Church, then, exists. From it (as present, as contemporary with the questioner, whereby the problem is upheld in the equality of contemporaneity) one can learn what is essentially Christian, for that is what the Church confesses. Correct. But at this point the matter cannot hold, in a historical-objective sense. For after it has been said of the Church, that it exists, and that from it one can learn what the Christian is, it is said again of this Church, the present one, that it is the apostolic Church, that it is the same Church that has existed for eighteen centuries. Now proof is required again. The only historical thing that is higher than proof is contemporary existence; every determination of the past requires proof. If someone were to say to another: prove that you exist, then
% --- p. 14 ---
the other answers quite rightly: that is nonsense. If, on the other hand, he says: I, who now exist, existed more than four hundred years ago, essentially the same, then the first rightly says: here a proof is needed. If a historical proof is now to be given for the apostolic origin and unaltered content of the confession of faith, it becomes, again, an approximating, and this has ``the remarkable property, that it can go on, for as long as it must.''
% Printer's defect (per transcription.tex's quote-balance ledger): the
% quotation opened before "cuts away," above, is never closed in the
% Danish -- no closing mark appears anywhere in the long stretch that
% follows, though the content plainly shifts to Nielsen's own commentary
% partway through (at "Correct."). Reproduced here with the same
% unclosed quotation, mirroring the print.

The proof of centuries for the truth of Christianity is unreliable, since it has, in its very first word, transformed itself into a hypothesis. ``A hypothesis can become more probable by persisting for three thousand years, but it never, therefore, becomes any eternal truth, that can be decisive for one's eternal blessedness. Has not Mohammedanism persisted for twelve hundred years? Eighteen centuries have no greater proving-force than a single day in relation to an eternal blessedness; eighteen centuries, on the other hand, and everything, everything, everything that can be told and said and repeated on that occasion, have a power of distraction, that distracts admirably.''

But, one might ask, would the pseudonym then perhaps have the critical introduction, together with the learned philology, done away with entirely? --- ``No, learned philology is entirely within its rights, and the present author, indeed, cherishes, despite everything, a reverence for what science holds sacred. But of the learned critical theology, on the other hand, one gets no clean impression. Its whole endeavor suffers from a certain conscious or unconscious duplicity. It constantly looks, as though something for faith, something concerning it, should suddenly result from this criticism. Therein lies the impropriety.'' Where, then, does the critical introduction lead? ---
% --- p. 15 ---
``If one has something up one's sleeve, then let us bring it out, to think it through in all dialectical calm. The one who, for the sake of faith, champions the Bible, must, indeed, have made clear to himself, whether anything in this respect would result from his whole labor, should it succeed beyond all possible expectation, so that he does not remain sitting in the parenthesis of his labor and forget, over the learned difficulties, the decisive dialectical closing bracket. The one who attacks must, likewise, have made clear to himself, if the attack now succeeded on the largest possible scale, whether anything else would result than the philological result, or, at most, a victory won by arguing \textit{e concessis}, where, mark well, one can lose everything in another way, if the mutual agreement is a phantom.''

``So that the dialectical may be done justice, and, undisturbed, simply think its thoughts through, let us assume, then, first the one thing and then the other.''

``So, I assume, that it has succeeded in proving of the Bible, whatever any learned theologian, in his happiest moment, could ever have wished to prove of the Bible. These books belong to the canon, no others, they are authentic, they are whole, their authors are trustworthy --- one can well say, it is as though every letter were inspired (more than this one cannot say; for inspiration is faith's object,
% --- p. 16 ---
is qualitatively dialectical, not to be reached by any quantifying). Furthermore, there is not a trace of contradiction in the holy books. For let us be hypothetically cautious: let there be so much as a hint of any such thing, and the parenthesis is there again, and the philological-critical restlessness leads one at once astray --- So, everything assumed to be in order with respect to the Holy Scripture --- what then? Has the one who did not have faith come a single step nearer to faith? No, not a single one. For faith does not result from a scientific deliberation, nor does it come about directly either; on the contrary, one loses, in this objectivity, the infinite personal interestedness in passion, which is faith's condition, the \textit{ubique et nusquam} in which faith can come to be. Has the one who had faith gained anything with respect to faith's power and strength? No, not the very least; rather he is, in this extensive knowledge, in this certainty that lies at faith's door and covets it, so perilously placed, that he will need much exertion, much fear and trembling, in order not to fall into temptation and confuse knowledge with faith. Whereas faith has, hitherto, had a beneficial disciplinarian in uncertainty, it would, in certainty, gain its most dangerous enemy. For if passion is taken away, faith no longer exists, and certainty and passion are not held in tension. What good fortune, then, that this wishful hypothesis, critical theology's most beautiful wish, is an impossibility, because even the most perfect realization would, after all, only become an approximation.
% --- p. 17 ---
And again what good fortune for the men of science, that the fault lies with them in no way at all! Were all the angels to join forces, they could still only produce an approximation, because, in relation to a historical knowledge, an approximation is the only certainty --- but also too little to build an eternal blessedness upon.''

``So, now, I assume the opposite, that the enemies have succeeded in proving of the Scripture, what they desire, so surely, that it surpasses the fiercest enemy's most ardent wish, --- what then? Has the enemy thereby abolished Christianity? By no means. Has he harmed the believer? By no means, not in the very least. Has he earned the right to exempt himself from the responsibility of not being a believer? By no means. For, since these books are not by these authors, are not authentic, are not \textit{integri}, are not inspired (this, however, cannot be disproved, since it is faith's object), it does not, therefore, follow, that these authors have not existed, and, above all, not, that Christ has not existed. To that extent it still stands just as free to the believer to assume it, just as free, let us note this well; for if he assumed it by virtue of some proof, he would be, in effect, giving up faith. Should it ever come to that, the believer would always bear some guilt, if he himself had invited it, and had begun by placing the victory in unbelief's hand, by wishing himself to prove it.''

$\beta$) \emph{The speculative way.} --- Speculation concerns itself only with Christianity as a given phenomenon,
% --- p. 18 ---
indifferent to whether there exist real, individual Christians or not; it assumes, without further ado, that we are all, roughly speaking, Christians. It is spoken of impersonally: ``Speculation does everything, speculation doubts everything, and so on. The one speculating does not speak of himself, for, in relation to speculation, one's own personal self is a matter of indifference. But the impersonal relation to Christianity is a misrelation. If Christianity is a matter of blessedness, then it is, indeed, only ``two classes of human beings who can know anything of it: those who, infinitely interested in passion for their eternal blessedness, believingly build it upon their believing relation to it; and those who, in the opposite passion (but in passion), reject it --- the happy and the unhappy lovers. Objective indifference learns absolutely nothing.''
% Printer's defect (per transcription.tex's quote-balance ledger): a
% second quotation opens here ("two classes...") before the first one
% ("Speculation does everything...misrelation.") is closed -- the Danish
% never supplies a closing mark for the outer quotation before nesting
% the inner one. Reproduced here with the same unclosed outer quotation.

Speculation regards Christianity as a historical phenomenon, and believes it can know its inward being from its outward form. That is impossible. In Christianity, the outward, as outward, is a matter of indifference and vanishing. ``Take a married couple; behold, their marriage plainly impresses itself upon the outward; it forms a phenomenon in existence (in miniature, just as Christianity has, world-historically, impressed itself upon the whole of life); but their marital love is no historical phenomenon; the phenomenal is the insignificant, has significance only for the married couple through their love; but seen otherwise (i.e., objectively) the phenomenal is a deception. So too with Christianity. The invisible Church is no historical phenomenon;

it does not let itself be regarded objectively as such at all, for it exists only in subjectivity.''

For the speculating one, the question of his personal eternal blessedness cannot, then, come forward at all, precisely because his task consists in coming more and more away from himself and becoming objective, and thus vanishing for himself and becoming ``speculation's beholding power.'' Since the human being is a synthesis of the temporal and the eternal, the blessedness of speculation, that the speculating one can have, will be an illusion, because he, in time, wishes to be merely eternal, wishes to be \textit{sub specie æterni}. Herein lies the speculating one's untruth. The infinite interest, in passion, for one's personal eternal blessedness is thus higher than the speculative blessedness, higher, because it is truer, because it definitely expresses the synthesis.

``In relation to an observation, where it is of importance, that the observer find himself in a definite state, it holds, indeed, that, when he is not in this state, he cognizes nothing at all. He can, now, deceive someone by saying, that he is in that state, although he is not; but, when it happens so fortunately, that he himself says, that he is not in the requisite state, then he deceives no one. If, now, Christianity is essentially something objective, then it holds for the observer to be objective; but, if Christianity is essentially subjectivity, then it is a mistake, if the observer is objective. In relation to all cognition, where it holds, that the object of the cognition is subjectivity's own inwardness, it holds, that the one cognizing must be in this state.
% --- p. 20 ---
But the expression for subjectivity's utmost exertion is the infinitely passionate interest in one's eternal blessedness.''

``The way of objective reflection makes the subject into the accidental, and thereby existence into a matter of indifference, a vanishing. Away from the subject goes the way to objective truth, and while the subject and subjectivity become a matter of indifference, so does the objective truth, and this is its objective validity, for the interest is, like the decision, subjectivity.
% Printer's defect (per transcription.tex's quote-balance ledger): this
% quotation, opened above, is never closed -- the paragraph ends without a
% closing mark, and a fresh quotation opens immediately with the next
% paragraph. Reproduced here with the same unclosed quotation.

``The way of objective reflection leads now to abstract thought, to mathematics, to historical knowledge of various kinds; it constantly leads away from the subject, whose existence or non-existence becomes, objectively, quite rightly, infinitely indifferent, quite rightly, for existence and non-existence have, as Hamlet says, only subjective significance. At its maximum this way will lead to a contradiction, and insofar as the subject does not become entirely indifferent to himself, this is, indeed, only a sign, that his objective striving is not objective enough; at its maximum this way will lead to a contradiction, that only objectivity has come to be, and subjectivity has gone out, that is to say: the existing subjectivity, that has made an attempt to become what one calls, in the abstract sense, subjectivity, the abstract form of abstract objectivity. And yet the objectivity that has come to be is,
% --- p. 21 ---
subjectively regarded, at its maximum, either a hypothesis or an approximation, because all eternal decision lies, precisely, in subjectivity.''

\begin{center}
\emph{b) Truth, inwardness, is subjectivity.}
\end{center}

``One believes, generally, that being subjective is no art. Well, of course, every human being is, indeed, also, in that sense, a piece of a subject. But now to become what one, in that sense, without further ado, already is: well, who would waste his time on that, that would be, indeed, the most resigned of all tasks in life. While we are all, in that sense, what one calls subjects, and labor to become objective, poetry goes about anxiously seeking its object; and yet poetry must, precisely, have subjectivities. Why, then, does it not simply take the first one at hand from our esteemed midst? Ah, no, he is no good, and if he will do nothing but become objective, then he will never become good for anything. This does, indeed, really seem to indicate, that there is something peculiar about being a subject. Why have a few become immortal as inspired lovers, a few as magnanimous heroes, and so on, if everyone, in every generation, were, in that sense, without further ado, already that, by, in that sense, without further ado, being a subject? And yet being a lover, a hero, and so on, is precisely reserved to subjectivity, for objectively one does not become it. Poetry would lift us up in the contemplation of the excellent --- of what excellent thing? Why, of subjectivity. So it is, then, something
% --- p. 22 ---
excellent, to be subjectivity. And now the priests! Why is it only a certain sum of pious men and women, to whose venerable remembrance the discourse always returns, why does the priest not simply take the first one at hand from our esteemed midst and make him a model: we are, indeed, all, in that sense, what one calls subjects. And yet piety lies, precisely, in subjectivity, objectively one does not become pious. The priests care for souls by carrying us away into devotion, while the soul's longing seeks after the transfigured --- what transfigured ones? Why, those who had faith. But faith lies, indeed, in subjectivity, so it is, then, something excellent, to be subjectivity.''

``It must constantly be maintained, that the subjective problem is not something about the matter, but is subjectivity itself. It is a matter of there being, objectively, absolutely no trace of any matter, for, in the same moment, subjectivity would slip away from some of the pain and crisis of decision, that is, would make the problem a little objective.''

``The objective way, however, believes itself to have a security, that the subjective way does not have (and it goes without saying, that existence, the act of existing, and objective truth cannot be thought together); it believes itself to avoid a danger, that awaits the subjective way, and this danger, at its maximum, is madness. By the merely subjective determination of truth, madness and truth become, in the end, indistinguishable, because they can both have inwardness. Not even this is true, however, for the inwardness of infinity never has
% --- p. 23 ---
madness; its fixed idea is, precisely, a kind of objective something, and madness's contradiction is precisely to want to embrace it with passion. What tips the scale with respect to madness is, then, again, not the subjective, but the little finitude, that becomes fixed, which infinity, indeed, can never become. Or can one not also go mad by becoming objective? The absence of inwardness is, indeed, also madness. Objective truth, as such, in no way decides, that the one who utters it is sane; on the contrary, it can even betray, that the man is mad, even though what he says is altogether true, and especially objectively true.''

Could the existing one really be outside himself, then truth would be something concluded for him. But where is this point? ``The I-I is a mathematical point, which does not exist at all; to that extent, anyone may well take up this standpoint, the one does not stand in the other's way. Only momentarily can the single individual, existing, be in a unity of infinity and finitude, that lies beyond the act of existing. This moment is the moment of passion. Modern speculation has driven everything to excess, so that the individual can objectively come outside himself, but this cannot be done at all. Existence holds one back. Modern speculation makes light of passion; and yet passion, for an existing one, is precisely existence's highest, --- and existing beings we are, after all. In passion, the existing subject is made infinite in the eternity of fantasy, and
% --- p. 24 ---
yet, at the same time, precisely most determinately himself. The fantastic I-I is not infinity and finitude in identity, for neither the one nor the other is real; it is a fantastic agreement in the clouds, a fruitless embrace, and the single I's relation to this airy vision is never stated.''

``The difference in way between objective and subjective cognition is now to be expressed thus:''

``When the question of truth is asked objectively, one reflects objectively upon truth as an object, to which the one cognizing relates himself. One does not reflect upon the relation, but upon its being the truth, the true, to which he relates himself. When that, to which he relates himself, is merely the truth, the true, then the subject is in the truth.''

``When the question of truth is asked subjectively, one reflects subjectively upon the individual's relation; if only the \emph{how} of this relation is in truth, then the individual is in truth, even if he thus related himself to untruth.''

``Let us take as an example: the cognition of God. Objectively, one reflects upon its being the true God; subjectively, upon the individual's relating himself to a something, \emph{in such a way}, that his relation, in truth, is a God-relation. On which of the two sides, now, is the truth? Mediation does not help. An existing one cannot be in two places at once, cannot be subject-object. When he comes nearest to being in two places at
% --- p. 25 ---
once, he is in passion, but passion is only momentary, and passion is precisely subjectivity's highest. The existing one, who chooses the objective way, now enters upon the whole approximating deliberation, that would objectively bring God forth, which is, throughout all eternity, never attained, because God is subject and, therefore, only for subjectivity in inwardness. The existing one, who chooses the subjective way, grasps, in the same moment, the whole dialectical difficulty involved in his needing to spend some time, perhaps a long time, in order to find God objectively; he grasps this dialectical difficulty in the whole of its pain, because he must have God in the same moment, because every moment is wasted, in which he does not have God. In the same moment he has God, not in virtue of any objective deliberation, but in virtue of inwardness's infinite passion. The objective one is not troubled by such dialectical difficulties as this: what does it mean, to devote an entire research-lifetime to finding God.''

``If someone, who lives in the midst of Christendom, goes up into the house of God, into the true house of God, with the true conception of God in knowledge, and now prays, but prays in untruth, and if someone lives in an idolatrous land, but prays with the whole passion of infinity, though his eye rests upon an image of an idol: where, then, is there most truth? The one prays in truth to God, though he worships an idol; the other prays in untruth to the true God, and thus worships, in truth, an idol.''
% --- p. 26 ---
``When one inquires objectively after immortality, another stakes the passion of infinity upon uncertainty: where, then, is there most truth, and who has most certainty? The one has, once and for all, entered upon an approximating, that never ends, for immortality's certainty lies, indeed, precisely in subjectivity; the other is immortal, and struggles, precisely for that reason, by contending with uncertainty. The one helps himself with a few proofs, and all his interest falls upon the proofs, for he is, indeed, objective, and objective knowledge is a matter of indifference to the existing individual; the other stakes his whole life upon an ``if there is an immortality,'' and the little bit of uncertainty helps him, because he himself contributes, with the passion of infinity.''

``When subjectivity is the truth, the determination of truth must, at the same time, contain within itself an expression for the opposition to objectivity, and this expression indicates, then, also, the tension of inwardness. Here is such a definition of truth: \emph{the objective uncertainty, held fast in the appropriation of the most passionate inwardness, is the truth, the highest truth there is for an existing one.}''

``But the given determination of truth is a paraphrase of faith. Without risk, no faith. Faith is precisely the contradiction between inwardness's infinite passion and objective uncertainty.''

\begin{center}
\emph{c) For the existing subjectivity, the highest truth
is the Paradox.}
\end{center}

The Paradox! Yes, here is the point, the ticklish point of contention, the problem, the Gordian knot, that ``science'' cannot untie. If a good deal of what has recently been said about the true and untrue cognition of God must already awaken some misgiving, among other things, because it is said so boldly and unconditionally, then surely the claim about the Paradox must awaken still greater offense, and be felt, involuntarily, as a thorn in the eye. I confess, that nothing else I have heard or read about the cognition of God has, to such a degree, both attracted and repelled me, as the bold claim, that the highest truth is a paradox. I admit, without reservation, that the pseudonym, on this point, more often than not seemed to me to work far more by a piquant whim than by any better insight; indeed, he sometimes even struck me as a perverse, seared thinker, bent on confusing everything, as a whimsical, cunning human being, who, by a \textit{coup de hasard}, would assassinate ``Christian science.'' But when, in such a fit of ill humor, I had thrown the book from me, and would have nothing more at all to do with the paradoxical fellow; he returned again, in all meekness, and began to dialecticize: ``Why did you break off the inquiry?'' --- It tears down truth. --- ``Which truth?'' --- Religion's. For the believer must, surely, be just as passionately interested in knowing, that it really is the true God he worships, as that he worships God in truth. ---
% --- p. 28 ---
``But can any human being at all come to a true cognition of God, by setting himself in an untrue relation to God? Do not be hasty; let us think the thoughts through once more, calmly! The true relation to God is, after all, not the objective one, in which God becomes a matter of indifference to the individual, as the individual becomes a matter of indifference to God; but it is the subjective one, it is the relation of faith and worship. When, now, the believer, infinitely anxious to gain certainty about the objective truth (that is, God's truth in himself), dares not step outside the subjective relation, so as not to fall into untruth; when he, ceaselessly laboring to attain the objective, only attains it, not together with the subjective, but as the subjective: is not this his endeavor, then, paradoxical?'' --- That is undeniable. --- ``Why, then, did you break off the inquiry?'' --- Because it offended my subjectivity. --- ``Yes, by way of subjectivity.''

It might, perhaps, seem strange, that I, on this occasion, thus expose myself and almost dramatize my subjective reflections. To this I can only answer, that it happens by no means because I now so gladly wish to be publicly put on display, or to make myself important through my private difficulties; but --- it has simply gone with me thus, with this ``Johannes Climacus,'' and I say it straightforwardly, because I suppose, that it might, perhaps, go with others, as it has gone with me.

So then: truth is the Paradox. ``When subjectivity, inwardness, is the truth, then truth,
% --- p. 29 ---
objectively determined, is the Paradox, and the fact that, objectively, truth is the Paradox shows, precisely, that subjectivity is the truth, since objectivity, indeed, repels, and objectivity's repulsion, or the expression for objectivity's repulsion, is the tension and measure of strength of inwardness.''

``Socratic ignorance is the expression for objective uncertainty; the existing one's inwardness is the truth. Socratic ignorance is an analogue to the determination of the Absurd, except that, in the repulsion of the Absurd, there is even less objective certainty, since there is only certainty that it is absurd, and, precisely for that reason, an infinitely greater tension in inwardness; Socratic inwardness is an analogue to faith, except that faith's inwardness, corresponding not to the repulsion of ignorance, but to that of the Absurd, is infinitely deeper. Socratically, the eternal, essential truth is by no means paradoxical in itself, but only by relating itself to an existing one. This is expressed in the Socratic proposition, that all cognition is a recollection. This proposition is an intimation of the beginning of speculation; but Socrates did not, therefore, pursue it further either; it essentially became Platonic.''

``The infinitely meritorious thing about the Socratic was, precisely, to accentuate, that the one cognizing is existing, and that existing is the essential thing. To go further by not understanding this is only a middling merit. This, then, we must, accordingly, have \textit{in mente}, and then see, whether the formula could not, after all, be so altered, that one really goes further than the Socratic.''

``So, subjectivity, inwardness, is the truth; is there a more inward expression for this? Yes, if the statement: subjectivity, inwardness, is the truth, begins thus: subjectivity is untruth. Let us not be hasty. Speculation, too, says, that subjectivity is untruth, but says it in exactly the opposite direction, namely toward objectivity's being the truth. Speculation determines subjectivity negatively, in relation to objectivity. The other determination, by contrast, hinders itself, in that it would begin, which precisely makes inwardness far more inward. Socratically, subjectivity is untruth, if it will not grasp, that subjectivity is the truth, but, for instance, will be objective. Here, on the other hand, subjectivity, in that it would begin to become the truth by becoming subjective, is in the difficulty, that it is untruth. So the labor goes backward, backward, that is, into inwardness. It is so far from being the case, that the way should thus lead toward the objective, that the beginning lies still deeper in subjectivity.''

``But untruth the subject cannot be eternally, or be presupposed eternally to have been; it must have become it in time, or become it in time.''

``The Socratic Paradox lay in this, that the eternal truth relates itself to an existing one, but now existence has, a second time, marked the existing one; such a change has taken place in him, that he can, by no means, Socratically, by recollection, take himself back into eternity.
% --- p. 31 ---
To do this is to speculate; to be able to do this, but to cancel the possibility of it by grasping inwardness in existence, is the Socratic; but now the difficulty is this, that what accompanied Socrates as a canceled possibility has now become an impossibility.''

``The Paradox arises, when the eternal truth and the act of existing are put together; but each time the act of existing is felt, the Paradox becomes clearer and clearer. Socratically regarded, the one cognizing was an existing one; now the existing one is so marked, that existence has effected an essential change in him.''

``Let us now call the individual's untruth \emph{sin}. Regarded eternally, he cannot be sin, or be presupposed eternally to have been in it. So, by coming to be (for the beginning was, indeed, that subjectivity is untruth), he becomes a sinner. He is not born as a sinner, in the sense, that he is presupposed as a sinner before he was born, but he is born in sin and as a sinner. This we may, indeed, call \emph{original sin}. But if existence has thus gained power over him, then he is prevented, by recollection, from taking himself back into eternity. If it was already paradoxical, that the eternal truth relates itself to the existing one; now it is absolutely paradoxical, that it relates itself to such an existing one. But the more difficult it is made for him to take himself, by recollection, out of existence, the more inward his existing can become within existence.''

``Let us now go further, let us assume, that the
% --- p. 32 ---
eternal, essential truth is itself the Paradox. How does the Paradox arise? By the eternal, essential truth and the act of existing being put together. When we, then, thus put it together in truth itself, truth becomes a paradox. The eternal truth has come to be in time, this is the Paradox. If the subject was, in what has just preceded, prevented from taking himself back into eternity by sin: now it need not trouble itself over this, for now the eternal, essential truth is not behind it, but has come before it, by itself existing or having existed, so that, if the individual does not, existing, in existence, take hold of the truth, he will never get it.''

``More sharply can existence never be accentuated, than it has now become. Speculation's deceit, of wishing to recollect itself out of existence, is made impossible. Here there can only be talk of grasping this; every speculation, that would be speculation, shows \textit{eo ipso}, that it has not grasped this. The individual can push all this away from himself and take refuge in speculation; but to accept it and then wish to cancel it by speculation is impossible, because it is calculated precisely to prevent speculation.''

``When Socrates believed, that God exists, he held fast the objective uncertainty with the whole passion of inwardness, and in this contradiction, in this risk, is precisely faith. Now it is different: instead of objective uncertainty, there is here the certainty, that it is, objectively regarded, the Absurd, and this Absurd, held fast in the passion of inwardness, is faith. Socratic ignorance is like a witty jest, in comparison with the earnestness of the Absurd,
% --- p. 33 ---
and Socratic existing inwardness like a Greek carefreeness, in comparison with the exertion of faith.''

``What, now, is the Absurd? The Absurd is, that eternal blessedness has come to be in time, that God has come to be, has been born, has grown, and so on, has come to be altogether as a single human being, not to be distinguished from another human being; for all immediate recognizability is pre-Socratic paganism, and, Jewishly regarded, idolatry; and every determination of that which really goes further than the Socratic must essentially bear the mark of God's having come to be, because faith, \textit{sensu strictissimo}, refers itself to becoming.''

The Absurd is, precisely, by its objective repulsion, faith's measure of strength in inwardness. The Absurd is, precisely, faith's object and the only thing that lets itself be believed. Insofar as the Absurd has the moment of becoming within it, the way of approximation will also be one that confuses the absurd fact of becoming, which is faith's object, with a simple historical fact, and thus seeks historical certainty for that which is precisely the Absurd, because it contains the contradiction, that what can only, directly against all human understanding, become the historical, has become it.''

``Christianity has now itself proclaimed itself to be the eternal, essential truth, that has come to be in time; it has proclaimed itself as the Paradox, and demanded the inwardness of faith
% --- p. 34 ---
in relation to what is, to the Jews, an offense, and to the Greeks, a folly --- and to the understanding, the Absurd.''

\begin{center}
\emph{d) In opposition to speculation's objective
knowledge, Christianity is to be designated as a
communication of existence.}
\end{center}

``When the several spheres are not held decisively apart from one another, everything becomes confused. When one is, thus, curious in relation to a thinker's actuality, finds it interesting, that one knows something about it, and so on, then one is intellectually to be blamed, because, in the sphere of intellectuality, it is precisely the maximum, that the thinker's actuality is a matter of complete indifference. But by being thus fussy in the sphere of intellectuality, one acquires a confusing resemblance to a believer. A believer is precisely infinitely interested in another's actuality. This is, for faith, the decisive thing, and this interestedness is not some slight curiosity, but the absolute dependence upon faith's object. The problem is expressed thus:

\emph{The individual's eternal blessedness is decided in time by the relation to something historical, which, moreover, is so historical, that there is taken up into its composition that which, according to its essence, cannot become historical, and must, therefore, become so in virtue of the Absurd.}

``Faith's object is not a doctrine, for then the relation would be intellectual, and it would be a matter, not of failing, but of attaining the maximum of the intellectual relation.
% --- p. 35 ---
Faith's object is not a teacher who has a doctrine; for when a teacher has a doctrine, the doctrine is \textit{eo ipso} more important than the teacher, and the relation intellectual, where it is a matter of not failing, but of attaining the maximum of the intellectual relation. But faith's object is the teacher's actuality, the fact that the teacher really exists. Faith's answer is, therefore, absolutely either yes or no. For faith's answer is not in relation to a doctrine, whether it is true or not, not in relation to a teacher, whether his doctrine is true or not, but is the answer to the question of a fact: do you assume, that he has really existed? And, mark well, the answer is given with infinite passion. --- Faith's object is God's actuality in the sense of existence. But to exist means, first and foremost, to be a single individual, and this is why thought must look away from existence, because the individual does not let itself be thought, but only the universal. Faith's object is, then, God's actuality in existence, that is, as a single individual, that is, that God has existed as a single human being.''

``Christianity is no doctrine of the unity of the divine and the human, of the subject-object, not to mention the other logical paraphrases of Christianity. For if Christianity were a doctrine, then the relation to it would not be one of faith, for to a doctrine there is only an intellectual relation. Christianity is, therefore, no doctrine, but the fact, that God has existed as a single individual on earth.''
% --- p. 36 ---
``Faith is, thus, no remedial lesson in the sphere of intellectuality, no asylum for weak minds. But faith is a sphere unto itself, and every misunderstanding at once recognizable by this, that it transforms it into a doctrine, and draws it within the compass of intellectuality. What counts, in the sphere of intellectuality, as the maximum, namely to become entirely indifferent to the teacher's actuality, holds, conversely, in the sphere of faith; its maximum is the \textit{quam maxime} infinite interestedness in the teacher's actuality.''

``If Christianity were a doctrine, it would \textit{eo ipso} not form an opposition to speculation, but would be a moment within it. Christianity indicates existence, the act of existing, but existence and the act of existing are precisely the opposition to speculation. The Eleatic doctrine, for instance, does not relate itself to the act of existing, but to speculation; a place is, therefore, to be assigned to it within speculation.''

``Precisely because Christianity is no doctrine, it holds, concerning it, that there is an immense difference between knowing what Christianity is, and being a Christian.''

``The question of what Christianity is must not be confused with the objective question of the truth of Christianity, which has already been treated above. Objectively, one can indeed ask, what Christianity is, in that the questioner would set this out objectively before himself, and let it stand undecided, for the time being, whether it is
% --- p. 37 ---
truth or not. (Truth is subjectivity.) The questioner then declines all reverend busyness with proving its truth, as well as all speculative haste to go further; he desires calm, desires neither recommendations nor haste, but only to learn, what it is.''

``Or can one not learn what Christianity is, without oneself being a Christian? All analogies seem to speak for it, that one can learn it, and Christianity itself must, indeed, regard as false Christians those who merely know what Christianity is. The matter has, here again, been confused by this, that one acquires the appearance of being a Christian, by being baptized straightaway as a child. But, when Christianity came into the world, or when it is introduced into a heathen land, it did not, and does not, simply draw a line through the contemporary generation of elders and take possession of the small children. At that time the relation was in order. Then it was difficult to become a Christian, and one did not concern oneself with comprehending Christianity; now we have almost reached the parody, that it is nothing at all to become a Christian, but difficult, and a busy task, to comprehend it. Hereby everything is turned about, Christianity transformed into a philosophical doctrine, where the difficulty, quite rightly, lies in comprehending, instead of Christianity's essentially relating itself to existence, and becoming a Christian being the difficult thing.''

``In relation to a doctrine, understanding is the
% --- p. 38 ---

maximum; and becoming a follower a cunning way, by which people who understand nothing sneak their way into acting as if they had understood: in relation to a communication of existence, existing in it is the maximum, and wishing to understand (comprehend) it a cunning evasion, that would shirk the task. To become a Hegelian is suspect, to understand Hegel is the maximum; to become a Christian is the maximum, to wish to understand (that is, comprehend) Christianity is suspect.''

``That one can know what Christianity is without being a Christian must, then, surely, be affirmed. It is another matter, whether one can know what it is to be a Christian, without being it, which must be denied.''

``What Christianity is must, then, first and foremost, be settled, before there can be any talk of mediating. This speculation does not enter into; it does not go about it in such a way, that it first sets forth what speculation is, then what Christianity is, in order then to see, whether the opposites let themselves be mediated; it does not first ascertain, to plain evidence, the respective identity of the contending parties in question, before it strives toward reconciliation. If one asks it, what Christianity is, it answers, without further ado: the speculative conception of Christianity, without troubling itself over whether there is anything in the distinction, that distinguishes between a something and the conception of a something, which seems here of importance for speculation itself; for, if Christianity itself is speculation's conception, then there is, indeed, no
% --- p. 39 ---
mediation, for then there is, indeed, no opposition, and a mediating between the identical is, after all, meaningless.''

``But even when speculation assumes a difference between Christianity and speculation, if for no other reason than to have the pleasure of mediating, when it does not, however, definitely and decisively state the difference; then one must ask: is not \emph{mediation} speculation's idea? When, then, a \emph{mediating} takes place between the opposites, the opposites (speculation --- Christianity) are not equal before the reconciler, but Christianity is a moment within speculation, and speculation gains the preponderance, because it had the preponderance, and because the moment of equilibrium, in which the opposites were weighed against one another, never arrived.''

``Yet the impropriety, of Christianity's being supposed a moment within speculation, has, surely, moved speculation to come somewhat to terms. Speculation has assumed the title: Christian, has wished to acknowledge Christianity by adding this adjective, just as sometimes, at marriages in noble families, a blended name is formed from both families', or as when trading houses unite into one firm, that nonetheless bears both their names. If it were now the case, as one so readily assumes, that becoming a Christian is nothing at all, then Christianity would, indeed, have to be overjoyed, that it had made such a good match, and had come to honor and dignity almost equal with speculation. If, on the other hand, becoming a Christian is
% --- p. 40 ---
the most difficult of all tasks, then the speculating one seems, rather, to profit, insofar as he wins, along with the firm, the being a Christian. But becoming a Christian really is the most difficult of all tasks, because the task, though the same, varies in relation to the respective individual's abilities. In relation, for instance, to comprehending, a good head has, straightforwardly, the advantage over the limited one, but not so in relation to believing. For when faith demands, that he give up his understanding, then it becomes just as difficult for the cleverest as for the most limited to believe, or, indeed, it perhaps becomes even more difficult for the former. One sees, again, the impropriety of transforming Christianity into a doctrine, where it is a matter of understanding, for thereby becoming a Christian comes to lie in the differential. What, then, is lacking here? The preliminary agreement, in which each party's status is settled, before the new firm is established.''

``Only let not, now, some quick head at once explain to a reading world, how foolish my whole book is, which one sees more than sufficiently from this, that I retail something like this, that Christianity is not a doctrine. Let us understand one another. One thing is, surely, a philosophical doctrine, that would be comprehended and speculatively understood; another thing, a doctrine, that would be realized in existence. When, in relation to this latter doctrine, there is to be talk of understanding, then this understanding must be to understand, that one is to exist in it, to understand, how difficult it
% --- p. 41 ---
is, to exist in it, what an immense existence-task this doctrine sets the learner. When, thus, at a given time, in relation to such a doctrine (a communication of existence), it has become common to assume, that being what the doctrine demands is very easy; but understanding the doctrine speculatively so difficult: then one can be in good agreement with this doctrine (the communication of existence), when one seeks to show, how difficult it is, existing, to comply with the doctrine. In relation to such a doctrine, it is, on the other hand, a misunderstanding to wish to speculate about it.''

``Such a doctrine is Christianity. To wish to speculate about it is a misunderstanding, and, as one goes further and further along this way, one incurs a greater and greater misunderstanding. If one finally reaches the point, of not merely wishing to speculate, but of having speculatively understood it, then one has reached the highest degree of misunderstanding. This point is reached by the mediation of Christianity and speculation, and, therefore, quite rightly, modern speculation is the highest misunderstanding of Christianity. Since this is now so, and since it is, moreover, given, that the nineteenth century is so terribly speculative, it is to be feared, that the word doctrine is instantly understood of a philosophical doctrine, that shall and will be comprehended. In order to avoid this impropriety, I have chosen to call Christianity a communication of existence, in order quite definitely to designate its dissimilarity from speculation.''

--- With this report, the review stands as a brief extract from an extensive treatment, a small loan from a great capital. The reader must now himself decide, whether he, with the reviewer, can here see a question of principle move, steadily and calmly, through a logically coherent train of thought, or whether the whole thing perhaps dissolves into ``scattered thoughts and aphorisms, whims and flashes.''

\begin{center}---\end{center}

After the review's foregoing, it will, surely, be discerned without difficulty, for what reason I now proceed directly to treat of the work recently published by Dr. H. Martensen: \emph{Christian Dogmatics.}

Since this dogmatic work will, presumably, be regarded as one of the most significant in more recent theology, and, as a modern scientific imprint of the old Lutheran type of doctrine, will be received with gratitude in the Danish Church, not only by the many candidates and younger clergymen, who owe their scientific formation chiefly to Dr. Martensen, but also by those who otherwise, with some ecclesiastical sense, take an interest in the conception of the Christian dogmas; it can, with justice, lay claim to a many-sided attention.

Were it now simply a matter of judging this extensive work with exclusive regard to the
% --- p. 43 ---
author's talent, learning, and gift of presentation, then I, for my part, would only have occasion to acknowledge once again the circumspection in the treatment of the material, the simple clarity of form, and purity of expression, that, in this work, as in ``Meister Eckhart,'' engage the reader from beginning to end. There is, indeed, surely, in this Christian Dogmatics, seen from the side of form, also something pleasing, something engaging, something that, in certain moods, works, if not strictly convincingly, then, at least, mildly soothingly. There is, in this Christian Dogmatics, something many-sided, something conciliatory, something peaceable, something accommodating. There is, in this Christian Dogmatics, when it is regarded as a systematic whole, also something distinctly artistic, something architectonic, something that at once, so to speak, weds the University building to the Palace chapel, science to religion: sharp understanding and naive simplicity, modern reflection and mystical immediacy, the clarity of the concept and the obscurity of the mystery: everything melts together in unity of style, in festive calm, solemn twilight, like the gleam of day mingled with the shine of an altar candle. It is a noble talent, a pregnant natural gift, whereby Dr. Martensen, as a university teacher, has inspirited so many disciples, and, as an ecclesiastical speaker, edified so many hearers.

As said, were it simply a matter of bringing out and acknowledging the beautiful formal side, then I would find, in the work here under discussion, passages enough that could serve me as evidence. Nor, either, will anyone who has taken note
% --- p. 44 ---
of my scientific endeavors readily be ignorant of how much regard I, in that respect, have for Dr. Martensen's talent and taste. But the question that is to be treated in this review is now no aesthetic question of form, but a decisive question of principle, and I therefore return to my first: \emph{Will Christianity, by its very nature, be an object of objective knowledge?} which question, in this context, can be held fast with the limitation: \emph{``Will Christianity be an object of dogmatizing speculation?''} --- This question, which has now become for me a kind of literary question of conscience, I shall, then, in this review, no longer lose sight of.

When Professor Martensen made his appearance at the University, and awakened, in the studying youth, the glad expectation, that Christianity would now let itself be explained within speculation, he had already (1837) published his first theological work: ``On the Autonomy of Human Self-Consciousness in the Dogmatic Theology of the Recent Age,''\footnote{Translated by L. B. Petersen.} and, in this work, indicated the presupposition for a Christian speculation. Under the same presupposition I then, likewise, worked out, some time afterward (1840), my own first treatise: ``The Application of the Speculative Method to Sacred History,''\footnote{Translated by Cand.\ theol.\ B. C. Bøggild.} a work of my youth, whereby
% --- p. 45 ---
I intended to prepare a corresponding presentation of the Life of Jesus. A presentation of the Life of Jesus I have, then, now also worked out,\footnote{\textit{The Gospel Faith and the Modern Consciousness, Lectures on the Life of Jesus.} Part 1, 1849, C. A. Reitzel.} but the execution does not correspond to the preparation: the standpoint has become another. By no means, that I have given up speculation, still less Christianity, but I have, always spurning the modifications of the middle-way doctrine, given up my earlier expectation of the fusion of the two, and have thus, on this point, changed my view, or, as it is called, shifted my standpoint.

It has gone otherwise with Dr. Martensen, who, as an organic genius, has remained true to his first love, and has not shifted his standpoint: his work on ``the Autonomy'' now corresponds to his Dogmatics, as the preparation to the execution, as the bud to the flower, the prophecy to the fulfillment.\footnote{``How I have let myself be taught by the signs of the times, with respect to the so much disputed question of the relation of cognition to faith, the work itself must show. To speak more fully of this in a preface I consider useless, since I have done so in the introduction, and everything, in the end, depends upon the execution. I may, however, be permitted to say, that, in the turn the philosophical movements have taken in the period last elapsed, I have found no occasion to give up the thought that underlies my first theological work: ``On the Autonomy of Human Self-Consciousness in the Dogmatic Theology of the Recent Age,'' the thought, that it is by no means religion that is to borrow its weight and
% --- footnote continues on p. 46 ---
significance from speculative thought, but that it is speculative thought that stands in need of religion, of God's revelation, as of its principle.'' That autonomy of self-consciousness, which I at that time (1837) sought to point out, has, since then, with astonishing rapidity, passed over into new forms of development.''\ldots{} From the preface to the \textit{Christian Dogmatics}.}
% --- p. 46 ---
Dr. M. has, indeed, in various respects, let himself be taught by ``the signs of the times''; but how this letting-himself-be-taught more precisely hangs together is provisionally indicated in the preface to the Dogmatics, and from this preface I set down, then, here the following:

``That a coherent theological thinking, indeed, theological speculation, is both possible and necessary, is a conviction I hope to be able to hold fast, even should it be demonstrated, that my own work had failed. Just as little as a negative speculation has been able to bring me to give it up, just as little have the objections been able to do so, which I have, of late, heard raised in the name of faith. Psychologically regarded, it must, indeed, be found altogether explicable, that there are, nowadays, many who distrust all science of faith, and think, that science belongs to free-thinking and paganism alone. At all times it holds true, too, that not science, not theology, but only faith, is the one thing needful for us all. Those, then, who feel no need to reflect upon their faith, may well be within their rights, when they, for their own part, regard cognition as unnecessary. And those who feel no drive toward coherent think-
% --- p. 47 ---

-ing, but can satisfy themselves by thinking in scattered thoughts and aphorisms, whims and flashes, may also be within their rights, when they, for their own part, regard coherent cognition as unnecessary. But, when it begins, of late, to be pronounced into a kind of dogma, that the believer can, absolutely, have no interest in seeking coherent cognition of that which is, to him, the highest; that the believer cannot wish to enter into any speculation, because all speculation is only cosmic, that is to say: worldly and pagan; that the believer must regard the concept of a science of faith as a self-contradiction, that abolishes genuine Christianity, and so on: then I confess, that such propositions, even when I have heard or seen them set forth with the paradox of wit, have had, for me, so little that is convincing, that I can see in them only a great misunderstanding, and a new, or rather, an old, misdirection. Thus much, however, I have been able to learn from it: that, just as there was a time, when we had not a few among us, who had become absolute in knowledge too soon, so, too, there is now, again, no lack of those, who have become absolute in faith too soon.''

Thus, then, has Dr. Martensen let himself be taught by ``the signs of the times''; so easily has he made short work of the objections, that have been raised, in the name of faith, against speculation; in this way he takes up the word on behalf of coherent thinking, and bears witness for objective knowledge.
% --- p. 48 ---
It is, indeed, very fine of Dr. Martensen, that he will still bear witness for ``Christian speculation,'' even at a time, when so many have fallen away, for this speaks, surely, for his having himself an inwardly firm and well-grounded conviction of speculation's necessity; it is, likewise, fine, that the dogmatician, who lives in a time, when unquiet thoughts all too easily disperse themselves into ``aphorisms, whims, and flashes,'' earnestly reminds us, how necessary it is, to seek a coherent cognition of that which ought to be, to the human being, the highest. We owe Dr. Martensen our thanks, both for the testimony and for the reminder; but how could we, indeed, thank him in a worthier way, than by striving, according to our slight ability, to make our way into his profound speculation, and, with all diligence, to search out, how matters now, in fact, stand with his own ``coherent cognition''?

But --- the task is difficult and the solution doubtful: empty assertions and stiff protestations prove, after all, nothing; it is, then, surely, best, to let the inquiry hover, as a case hovers before the court, and, meanwhile, to keep close to the question:

\begin{center}
{\large\bfseries Has not ``the Christian Dogmatics'' undialectically
evaded the Christian problem?\par}
\end{center}

However one may, otherwise, wish to determine the mutual relation of Christianity and speculation: this, surely, one must nonetheless admit, that they, each in its own form, at least,
% --- p. 49 ---
lay claim to expressing the Absolute. Christianity communicates the Absolute in the form of faith, speculation in the form of knowledge. The one, then, who holds to speculation, will, through knowledge, stand in an absolute relation to the Absolute, and, to that extent, be ``absolute in knowledge.'' The one, on the other hand, who holds to Christianity, will, through faith, stand in an absolute relation to the Absolute, and must, then, surely, to that extent, be ``absolute in faith.''

To declare oneself ``absolute in knowledge'' is a presumption, insofar as the speculating one thereby means to say of himself, that he has, in every respect, an absolute and perfect knowledge of the Absolute; but insofar as he thereby means to express, that the Absolute is, and must be, to him, the sole and sovereign principle in all his thinking and knowing, this is so far from being any presumption, that it is, much rather, a dutiful and necessary self-limitation. To be ``absolute in knowledge'' is, then, the same as to be strict, conscientious, and exacting with one's knowledge of the Absolute. The one who knows dare not expect any assistance from elsewhere, for instance from faith, so that the Absolute should confirm itself to him half in faith and half in knowledge; for, in that case, he would have violated his principle and corrupted his science, just as a mathematician has, indeed, violated his principle and corrupted his science, when, instead of stringent proofs, he argues from immediate experience. Even if the one who knows, as an existing individual, should need many other
% --- p. 50 ---
clarifications, and quite another conception of the true, than that which can be given him in knowledge's absolute form, nonetheless none of this concerns his knowledge. As one who knows, he has absolutely no other interest, and no other need, than that of coming to know the Absolute. Only on these terms can speculation develop, and vindicate its right, and make progress; only on these strict terms have the great thinkers, a Fichte, a Schelling, a Hegel, worked their way forward, and become, each according to his standpoint, ``absolute in knowledge.''

To wish to go about and pass for ``absolute in faith'' is a presumption, if the single individual thereby means to give himself out as able to accomplish the whole work of faith, to have, in every way, an absolute and unconditional faith, the faith that moves mountains; but, when it is, on the other hand, understood in such a way, that it is the same as making faith the absolute, unwavering condition for the existential appropriation of Christianity's Absolute; then this is so far from being a presumption, that it is even, straightforwardly, the believer's duty and obligation. If faith is the only right form and condition, under which the absolute truth of Christianity can be appropriated: then, indeed, the principle of faith, too, must be, in itself, an absolute principle, a principle that is only violated and corrupted, when it is to be propped up by something else, for instance by speculative knowledge. Even if the believer, as a thinking spirit (and then, in particular, as a genuine genius), may feel a need for much other knowledge, and quite another cognition,
% --- p. 51 ---
than that which is directly contained in the Absolute of faith: this, nonetheless, stands firm, that he, as a believer, neither has, nor can have, any other need, or other interest, or other task, than that of laboring for the appropriation, and accomplishing the work of faith. Only on these terms can the cognition of faith be clarified and furthered, and vindicate its right; only on these strict terms have the great witnesses of faith, from Paul all the way to Luther, suffered and struggled, and become, each according to his gift of the spirit, ``absolute in faith.''

Let no one suppose, that I, with these remarks, imagine myself to be saying anything new, or, indeed, anything that Dr. Martensen does not know. Far from it! Dr. M. knows it all just as well as I, and can, when he wishes, say it, moreover, far more beautifully. At any rate, he has, in the past, known how to speak so absolutely of the Absolute, that it was a true joy to listen to him. But what I say here, I say in perplexity, because it seems to me questionable, because it seems to me, that my most reverend friend, ``taught by the signs of the times,'' has, in an important matter, begun to change his language, without, however, having changed his standpoint. See, that is a circumstance that makes me perplexed, so that I hardly know what to say, a questionable turn, that makes my speech so equivocal, that I could even be tempted to consult ``the critical introduction,'' and seek counsel from ``hermeneutics.''
% --- p. 52 ---
How much is it, then, that Dr. M. has been able to learn from ``the signs of the times''? Answer: ``that, just as there was a time, when we had not a few among us, who had become absolute in knowledge too soon, so, too, there is now, again, no lack of those, who have become absolute in faith too soon.'' So, then: this much! at least, this much is said with definiteness; for what is said in what follows is nowhere near so definite. But how much, now, is this ``this much,'' when we look more closely? Well, that is precisely what I cannot make out. If I take the two expressions, ``absolute in knowledge'' and ``absolute in faith,'' in such a way, that they aim at the conceit that all too readily accompanies both speculative upswings and religious awakenings, then the meaning becomes this: ``when speculation reigned, we had not a few among us, who imagined themselves speculatively omniscient''; and if it is to be understood in this way, then Dr. Martensen is, indeed, quite right, that such persons must, surely, have come ``too soon'' with their knowledge. But the interpretation does not fit; for however much immaturity may, at the time, have shown itself among the speculating youth, it has, surely, nonetheless, been only very few (if there were any at all), who imagined themselves to have attained speculative omniscience. We surely owe Dr. M. enough respect, not to go and interpret his words in so unreasonable a way, that everything he has learned from ``the signs of the times'' dissolves into pure nonsense. This can, moreover, all the more easily be averted,
% --- p. 53 ---
in that the interpretation here proposed, which fits so poorly the first term (those absolute in knowledge), luckily does not let itself be applied to the second (those absolute in faith) at all; for those who, in this way, should have come ``too soon,'' would, indeed, have to be pure religious enthusiasts; but of such there is, surely, as yet, none among us.

But when I now understand the expressions, ``absolute in knowledge'' and ``absolute in faith,'' as referring to an absolute acknowledgment, both of the principle of knowledge and of that of faith, I am, still, no nearer. How can Dr. M. regard it as a new, and yet old, misdirection, that some should wish to be absolute in the principle of knowledge, others absolute in that of faith, when Dr. M. himself, indeed, according to his own standpoint, has as good as committed himself to being absolute in both of them?

To place the whole emphasis on ``too soon'' does not help; for to arrive too soon at an absolute acknowledgment of an absolute principle would be, indeed, the same as arriving too soon at the cognition of an essential truth. But it is, nonetheless, something strange, to complain that human beings arrive too soon at the cognition of truth; one is otherwise accustomed to complain, that they arrive too late. That Dr. M. must have learned something definite from ``the signs of the times,'' indeed, something quite particular, I am convinced of, since he says so himself; but what this something actually is, and what this something is actually meant to signify, I am not able to make out, unless it should be discoverable in what follows.
% --- p. 54 ---
``As far as I,'' --- Dr. M. continues --- ``am able to discern, there is only One who fully corresponds to the concept of the believer, namely the whole universal Church. Each one of us, as an individual, possesses faith only in a certain limited measure, and must, surely, guard ourselves against making our own individual, perhaps somewhat one-sided, perhaps even somewhat sickly, life of faith the rule for all believers. Only that great single one, who, as the Apostle teaches us, is, through the ages, to grow up into a perfect man, can fully realize the concept of the believer in the whole of its spirituality, the whole of its many-sidedness and rich multiplicity, can alone possess the fullness of faith and of all the gifts of grace. But of that great single one we know, by the Apostle's testimony, that he is in a continued advance toward the unity of faith and cognition. (Eph.\ 4:13).''

If it were not now Dr. Martensen, who spoke to us in these remarkable lines, I, for my part, should soon be done with the explanation; for profound seers, who interpret ``the signs of the times,'' without one's therefore being able to say, what it actually is, that they have seen; clever people, who, in disputes of principle, give good advice, without one's daring, on that account, to conclude, that they actually know, what the dispute is about; literary authorities, who, in refined superiority, warn society against certain individuals' excesses, without, therefore, finding it fitting to name these individuals by name: such things are, after all, nothing so very unusual, either in theology or in philosophy, and one
% --- p. 55 ---

need not take the matter to heart any further. But when a man who wills something, and knows what he wills, a man who weighs his words, and knows what he says, a man like the author of ``the Christian Dogmatics,'' who both can and will make a fundamental theological question clear to us, begins by speaking in riddles, then one grows uneasy in mind, unsure of oneself, and falls into a painful embarrassment. Should there really be found among us some single individuals who have, by some mistake, confused themselves with the ``holy, universal Church,'' some single individuals who fancy themselves to possess the perfect faith; some self-willed single individuals, who not only wish to break loose from the Church, but into the bargain make their own ``individual, one-sided, and morbid life of faith into a rule for all believers''? But how, then, does this hang together? Who are these single individuals? Where are they to be found? What is one to call them? --- I look and look --- at ``the signs of the times''; I see a numberless throng of individuals, of the most various minds and thoughts, but these particular single individuals I cannot see. ``That great single individual,'' of whom the Apostle speaks, I see, at a pinch, that is to say: I glimpse him as through a fog (for I had, otherwise, never before had the vision, that the universal Church should really be the single individual); but the small single individuals, of whom Dr. M. speaks, no, though I put on seven pairs of spectacles: these small, refractory single individuals I truly cannot see. Is there, then, a new heresy afoot? Does ``the single individual'' perhaps want the doctrine of the Church overturned, or the dogmas dis-
% --- p. 56 ---
-solved, or some point or other in the evangelical-Lutheran confession of faith distorted? Should the dogmatician ever happen upon such a heretic, he ought, assuredly, not fail to look him well over, and combat him thoroughly; for a teacher in the Church does well to take heed to himself and to the doctrine (1 Tim. 4:16). But if there is, now, no such heretic to be found? Well, then the dogmatician does well to take heed to himself, and let ``the single individuals'' be in peace.

What is it, then, that the Apostle Paul, according to Dr. M.'s ``hermeneutics,'' is really meant to teach us about that ``great single individual''? --- ``That he is in a continued advance toward the unity of faith and cognition.'' --- What cognition? For this can, surely, scarcely be the existential cognition of faith; unless one wishes to maintain, that Christendom, in the nineteenth century, has a far more inward and truer cognition of faith, than in the time of the Apostles. Far be it from me to wish to attribute to Dr. M. so unreasonable an opinion, especially since he, indeed (namely in §§ 185--189), himself maintains, in the most definite terms, the inspiration of the whole first period. Can it, now, not be the subjective cognition of faith, that is here spoken of, well, then it must, surely, be the objective, and \textit{in specie} the speculative. Should this be the meaning, then one can, without hesitation, grant the author to be right, that ``Christian speculation'' is, in the nineteenth century, far further advanced than in the first, and that we, to that extent, see herein ``a continued advance''; but what, then, follows from this?
% --- p. 57 ---
That ``the single individuals'' must not be too absolute in knowledge. And why not? Because those who become absolute in knowledge too soon do not give themselves time to see, that speculative knowledge must be supported by faith. But why, then, must the single individuals not be absolute in faith? No, because those who become absolute in faith too soon naturally forget, that faith must be supported by speculative knowledge. The single individual, that is to say: the little single individual, dare, then, no longer declare himself unconditionally either for faith or for knowledge, dare no longer, in any respect, speak absolutely of the Absolute, since ``that great single individual,'' according to the Apostle's, or rather according to Dr. Martensen's testimony, has not yet himself attained the Absolute, but ``finds himself in a continued advance toward the unity of faith and cognition.'' Therefore:

\begin{center}
\emph{a) ``The Christian Dogmatics'' gropes after the problem, but does not find it.}
\end{center}

That the question of the relation of faith to cognition stands in close connection with the question of the relation of dogmatics to philosophy, namely to the philosophy of religion, Dr. M. has seen very well; just as he has, in general, seen very well, which points, in so difficult an inquiry, might especially be brought forward. Straightway, in § 2 of the Introduction, he declares himself, then, also in the most definite terms for a separation of \emph{dogmatics and philosophy.}

``Dogmatics'' --- it says expressly --- ``is not merely a science about faith, but also a cognition
% --- p. 58 (NB: folio misprinted "53" in the original) ---
in faith and out of faith. It is not merely a cognition of what has been, or is, truth for others, without being so for the presenter himself; but neither is it a philosophical cognition of the Christian truth, which the presenter had won by taking up a standpoint outside faith and outside the Church. For even were we to assume --- which we, however, by no means grant --- that a scientific insight into the Christian truth is possible without Christian faith, such a philosophizing about Christianity, even were its results ever so favorable to the Church, could nonetheless not be called dogmatics. Dogmatics stands in Christendom, and the dogmatician is only the organ of his science, in that he is, at the same time, the organ of his Church, which does not hold of the mere philosopher, who will serve pure science alone.''

When I read this passage the first time, it seemed to me both set forth with clarity, definiteness, and beauty, and that I truly also understood, what it was, that was being said; but --- be that as it may --- when I read it a second time, I nonetheless did not understand it after all. The difficulty of the task can be expressed thus: there is delivered (as is required% <-- print "forbret" (d/b glyph); corrected
{} in § 1 of the Introduction) a \emph{scientific presentation and grounding of the Christian articles of faith in their coherence}; what significance, now, ought to be attached to such a presentation? Answer: if the author is a philosopher, who will serve pure science alone, then the presentation, even if its results are ever so favorable to the Church, does not
% --- p. 59 ---
therefore yet have the rank of dogmatics; but if the author, on the other hand, will append the oral or written declaration, that he has, indeed, worked out his ``scientific presentation'' most nearly, yes, most immediately, for the Church's sake, yes, that he ``is only the organ of his science, insofar as he is the organ of his Church'': well, then it does, indeed, with the help of the declaration, and with the Church's assistance, and with God's help, become dogmatics after all, for dogmatics stands in the Church, and --- the Church stands in the dogmatician. But what if ``the scientific presentation'' were now by an anonymous author; what if the anonymous author had known how to conceal himself so cunningly under the personification of ``that great single individual,'' that neither ``the critical introduction'' nor ``the hermeneutics'' were able, with any certainty, to determine, whether he really found himself within the Church or outside the Church: what then? Yet this is but a hypothesis; the passage itself is hypothetical; it will, indeed, \emph{by no means grant}, that a scientific insight into the Christian truth is possible without Christian faith. But how, then, are we to understand the connection, whereby what can by no means be granted on p. 4 can nonetheless, after all, be granted on p. 14?

``The religious view,'' --- it says in the note on p. 14 --- ``can also be possessed at second hand, that is to say, in an aesthetic or philosophical way. And precisely because the religious view moves in the element of objectivity, of thought and imagination, it can be loosed from its root of life in the disposition, and, independently of personal faith,
% --- p. 60 ---
be possessed in a merely aesthetic and philosophical way. Thus there are philosophers, poets, painters, sculptors, who, with great plastic power, have presented the Christian conceptions, without themselves, however, having possessed these conceptions in a religious way, in that they have only related themselves thereto through the medium of thought and imagination.''

It is truly a pity, that these two important passages (p. 4 and p. 14) hold themselves so widely apart in the book; had they stood nearer to one another, the author would, without doubt, have become aware of the apparent contradiction, and, for the reader's sake, have added a small clarifying intervening clause. As they now stand, they resemble two relations, who have parted in mutual disagreement, and turn their backs on one another. --- For once one has granted, that there are philosophers, who, without themselves being religious (i.e. without true Christian faith), have presented the Christian conceptions with great plastic power (that is, with philosophical, that is to say: scientific insight into the truth of these conceptions); well, then one might, indeed, if for no other reason, then at least for consistency's sake, just as readily grant, that a scientific insight into the Christian truth is possible without Christian faith; but this is now precisely the very thing, that one ``will by no means grant.''

That Dr. M. wills, in the most definite terms, to have dogmatics kept separate from the philosophy of religion, is clear enough, not, indeed, just from the execution, but nonetheless from the introduction. Not even Christian philosophy dares here coincide
% --- p. 61 ---
with dogmatic speculation. It says (p. 81, § 35), among other things, thus:

``But after Christianity itself has given birth to a new cognition of itself, a theology, the question arises, whether, alongside theology, there also remains room for a Christian philosophy, and how these, then, relate to one another. We presuppose here, that there is such a thing as a Christian philosophy; we presuppose, further, that it is bound to the same fundamental conditions of cognition as theology, and must, therefore, proceed from \textit{credo ut intelligam}; but we place the difference in this, that philosophy, even as Christian, is worldly wisdom, that is to say, cognition of the universe, whereas theology is cognition of God, a difference that is, indeed, relative, but a difference nonetheless.''

Yes, indeed, this difference between Christian philosophy and dogmatic speculation is a relative difference; it is by no means absolute, but relative, only relative: that is a true word, well said; but let us now examine, whether there really is, here too, any difference at all, or whether the relative difference does not, in the end, creep in and become smaller and smaller, whether it does not, in the end, become so small and so fine, that no one can produce it, and no one can get to see it, because, on account of its infinite relativity, it is, in the end, found to have altogether gone astray and vanished.

The difference is determined thus (p. 82): ``Philosophy seeks to cognize the divine law of the world, that
% --- p. 62 ---
accomplishes itself through the whole of existence, through the various spheres of the natural and the spiritual world, and it seeks to cognize Christianity as the fullness of the world's law. Philosophy, then, immerses itself in the manifold and leads this back to the One, to the Kingdom of God as the all-clarifying midpoint; theology, dogmatics, on the other hand, takes its standpoint from the outset in the center, immerses itself exclusively in the One, in the Kingdom of God as such.'' --- But, in that we, too, now immerse ourselves in what is here said about ``the One,'' about ``the Kingdom of God as such,'' we ought, surely, not to fail either to immerse ourselves in all that Dr. M., in so many other passages, has said about ``the cosmic miracle,'' about ``creation and cosmogony,'' about ``man and the angels,'' about ``sinful natural determination and sinful history,'' about ``the demonic powers and the Devil,'' about ``death and the perishability of all creation,'' about ``the free course of the world and God's manifold wisdom,'' about ``the Father's Logos,'' which is not merely ``the heart of God the Father,'' but at the same time the eternal ``heart of the world,'' etc., etc.; for all this points back to what is expressly said to us (p. 13) about religion and revelation, namely, that ``the idea of God takes shape in a comprehensive \emph{view} of the world and of human life in its relation to God, a view of heaven and earth, nature and history, heaven and hell.''

Is, then, the religious view really so comprehensive, and shall all this be thought through, investigated, and set forth with scientific, indeed with speculative-scientific, thoroughness; then
% --- p. 63 ---

it does, indeed, seem undeniable to follow from this, that so comprehensive a dogmatic cognition of God cannot possibly be separated from an equally comprehensive philosophical cognition of the world, i.e. a dialectical-speculative cognition of the universe. But what difference is there, then, between dogmatics and the philosophy of Christianity? Answer: none, none at all! The philosopher must take his standpoint in ``the center'' just as much as the dogmatician; for how should he cognize Christianity as ``the fullness of the world's law,'' without cognizing it as ``the central''; but how should he cognize the central, without taking his standpoint in the center? It is in vain, that the dogmatician appeals to faith and puffs himself up in the Church, when he% <-- print "har" (n/r glyph); corrected
{} says, speculatively: ποῦ στῶ,\footnote{Greek: ``where am I to stand''; the classical (Attic) interrogative form. The philosopher's answering πᾷ στῶ, two lines below, is the Doric form of the same question, and echoes Archimedes' saying, ``Give me a place to stand [δός μοι πᾷ στῶ], and I will move the earth.'' Nielsen's pun turns on the two dialect forms meaning, in effect, the same thing.} for the philosopher, too, now says, speculatively: πᾷ στῶ; and see, then they both stand, indeed, right in the middle of speculation. If the dogmatician, i.e. the speculating dogmatician, cannot dispense with ``the new categories of the world,'' then the philosopher of religion, too, cannot neglect ``the old categories of revelation''; for speculative productivity is precisely conditioned by the mutual development of both. ``Christian speculation,'' which, indeed, aims precisely at explaining the universe in Christianity, and, conversely, Christianity in the universe, cannot possibly separate these two functions, without dissolving and warping itself.

``But,'' --- one says --- ``there is, nonetheless, all the same a difference, namely in the treatment of the material there is a difference, a not insignificant difference; for, while the dogmatician, who has not merely πᾷ στῶ in speculation, but besides
% --- p. 64 ---
that a wholly different, higher, and far more reliable πᾷ στῶ in the Church, goes into detail and takes exact account of the several positive articles of doctrine; the philosopher, on the other hand, is accustomed only to treat the matter wholesale, and to pass over many ecclesiastical points, merely in order to get the philosophy of Christianity taken up as a link in his philosophical encyclopedia.'' --- But, in relation to the principle, this is, indeed, an altogether meaningless difference. For, if the speculating dogmatician takes up more material than speculation allows, he thereby, indeed, warps the science; and, conversely, if the philosopher neglects to appropriate all the material speculation requires, he likewise, indeed, warps the science: where, then, is the difference?

``Yes --- but --- yes,'' --- one says --- ``here is, nonetheless, all the same a difference, a certain difference, yes, indeed, here is a difference, it cannot be denied.'' --- Wherein, then, does the difference consist? --- ``Wherein it consists? Well, that cannot, of course, be said just like that with decisive definiteness, since sound reflection, indeed, at once shows, that the boundary we here have before us must be a fluid boundary, and, to that extent, also a boundary in the process of becoming. Sound reflection (cf. p. 80) will always recognize, that the speculative comprehending is itself a highly mobile and dialectical concept, that does not let itself be dismissed with a dry yes or no, does not let itself be dismissed with the assertion, that it must either be completely, or not be at all; for it is only as \emph{becoming}. So much can, nonetheless, meanwhile, be said with definiteness, that the dogmatician stands in the Church, and that he
% --- p. 65 ---
is only the organ of his science, in that he is, at the same time, the organ of his Church, which does not hold of the mere philosopher, who will serve pure science alone. If this, now, stands firm, that the dogmatician stands in the Church, and loves the faith as he loves the Church, then it ought also to be said with definiteness, that the intellectual love of the Christian truth, which is especially to be found in \emph{the teaching order of the Church}, is far higher, far more personal, and far more pathological, than that which is commonly found in the mere philosopher.'' --- But, when it is now well known, that this separation rests upon an unfortunate confusion% <-- print "Forverling" (r/x glyph); corrected
{} of the personal and the ingenious; when it is now well known, that one may be a believing Christian, and labor vigorously in faith, without having the intellectual, that is to say, the theoretical love, of which mention is here made, while another may be altogether governed and possessed by this objective love, without laboring, to the same degree, for the subjective appropriation in faith; when it is well known, --- what the author of ``the Christian Dogmatics'' assuredly knows as well as anyone --- that the intellectual and speculative love, with which the thinker immerses himself in the mysteries of revelation, is by no means any proper religious requirement (``faith is, indeed, the one thing needful for us all''), but solely and alone the drive and need of talent for the contemplation of the idea, for objective joy in divine visions: where, then, is the difference?

Assuredly, we cannot deny Dr. M. the acknowledgment, that he, as a dogmatic genius, has been intent
% --- p. 66 ---
upon making a definite separation between dogmatics and the philosophy of religion (a separation that, as is well known, is of an extraordinary importance); but all the more must we regret, that this talent has not yet succeeded in finding the point from which the separation proceeds. The examples he cites, in passing, are as unfortunate as can be conceived. In order to vindicate for dogmatics its independent principle, he points to Schleiermacher (p. 7); but what can that help, when Schleiermacher's standpoint (p. 11, § 7) is mistaken, and his dogmatic conception of religion is a, admittedly, very ingenious and profound, but nonetheless unsuccessful attempt, that neither agrees, nor can agree, with the orthodox doctrine of the Church? In order to express the speculative relation of objectivity to Christianity, the author refers (p. 76) to ``speculative mysticism and theosophy (Jacob Böhme), which teaches, that in faith itself there is a beholding''; but --- just as though it were, nonetheless, predetermined, that the reader should have neither profit nor joy of this beholding --- he again calls the reference back, and makes us aware, that here, too, there was a ``misdirection, that pointed away from history.'' How unfortunate, indeed how unpleasant, all this is! If dogmatics is to be separated with definiteness from the philosophy of religion, then the difference must, absolutely, be capable of being indicated in such a way, that these sciences are seen either to have an altogether different task, or, if the task is found, at a certain point, to be the same, then to solve it in so different a manner, that the methods, at the very
% --- p. 67 ---
least, diverge, as, indeed, the sides of an angle diverge, even though they begin from the same vertex. But --- unless I have unhappily blinded myself upon all the speculative splendor, that is poured out over the Martensenian introduction --- the intended separation ends in a sheer nought; for the standpoint, that is to say: the scientific standpoint (what the dogmatician, for the rest, does in the Church, does not concern this matter) remains the same, the task the same, the solution the same.

And Dr. M. will take up the word on behalf of ``coherent cognition''! It may, indeed, be all very well with coherent cognition and ``the fluid boundary,'' when everything else, moreover, will take shape, when dialectic works briskly and everywhere effects the requisite distinction; but when the differences flow together in the fluid, while the coherence condenses itself into a mystical half-darkness, that conceals the problem: then one ought, surely, not disdain a good ``aphorism,'' a ``scattered thought,'' or a ``flash,'' especially if it is a flash of lightning.

\begin{center}
\emph{b) ``The Christian Dogmatics'' will appropriate speculation, but without grasping speculation's problem.}
\end{center}

It cannot, assuredly, be denied, that Dr. M., indeed, in many ways and under many different turns of phrase, strives to impress upon the readers of the Dogmatics the indispensability of an objective knowledge. In this character he appears, both on behalf of the idea and on behalf of the Church, as spokesman for the higher science and namely for
% --- p. 68 ---
speculation. It is an honorable office, a handsome position, especially in a small country, where the number of those who feel a need and a calling to occupy themselves thoroughly with speculative philosophy cannot, ordinarily, be very great. Against this appearance there is, then, in and of itself, so far as I can judge, nothing to object; and I, for my part, would gladly, to the best of my ability, stand by Dr. M. in his estimable striving. But what somewhat alarms me in this case is the sharp sidelong glance, with which he regards certain ``single individuals,'' who, after first having tried themselves a little in speculation, and thereby having become, ``all too soon,'' ``absolute in knowledge,'' now, next, in more recent time, have cast speculation clean away, and have become, ``all too soon,'' ``absolute in faith.'' It is truly a vexing circumstance with such veiled utterances, that make one so unsure, so embarrassed, precisely when one has the very best will to serve the very best cause. One has an anxiety about offering one's service, when one does not rightly know how it will be received: perhaps one will be turned away as --- unusable, perhaps one has already been turned away; perhaps ``the Christian wisdom''\footnote{For ``the Christian wisdom'' --- yes, that I must confess --- ``the Christian wisdom'' knows how to make its point.} has already discarded one as an imitator, as --- \textit{miserabile dictu} --- a ``trivial imitator,'' and let one go off at the last reduction; perhaps one then, sometime, by chance, gets to learn, that there is
% --- p. 69 ---
one, who has been let go from his science, just like the man who, reading it in the newspaper, learned by chance, that it was he who had been dismissed from his office. It is a vexing circumstance; for however one twists and turns it, something mystical, nonetheless, mixes itself in here. I will, then, take the matter as plainly and simply as possible: is it really I, who have been judged, reduced, and discarded, deposed from ``coherent thinking'' and dismissed from ``objective knowledge'': well, then I shall have to resign myself to it, and it shall not matter, either, so long only as my sentence may be publicly declared to me, that I may get to know my fate. But should the veiled allusions perhaps not strike me quite dead; perhaps I can, then, provisionally, rescue myself out of the embarrassment, by adding a remark, concerning my philosophical striving. In the ten years, that I have now, thus, \textit{con amore et ex officio}, occupied myself with philosophy, the newer and the older, I have, with particular interest, followed speculation along its various ways, and thereby had occasion to see, how far, roughly, the hitherto discovered domains of metaphysics might, indeed, extend. Even though I myself have not discovered anything worth mentioning, I have, nonetheless, from my teachers, gained a strong impression of what ``coherent cognition'' has to signify, and --- what I regard as the most gratifying yield of my studies --- learned to admire the true heroes of systematic thinking, and to bow before the superiority of genius. Though I am, indeed, assuredly, no decided philosophical talent
% --- p. 70 ---
(still less a genius), I nonetheless credit myself with head enough, that I can, after all, tolerably distinguish speculation from fancy, and follow along fairly well, especially when a more gifted one will lead the way. So I let, then, Dr. M. lead the way again --- with the speculative theme, and follow him, as best I can, in order, nonetheless, to accompany him with a little dialectic.

Many and remarkable are the passages, that, in ``the Christian Dogmatics,'' speak to the exaltation of speculation. ``For,'' --- it says (p. 78) --- ``though we fully grant, what Melanchthon\ldots{} says against an idle speculation, that was loosed from life, the cognition of God's mysteries is nonetheless, in itself, a good, and the cognition of God's glory is, in itself, edifying.'' Further (p. 79): ``The thorough explication cannot but develop such thought-oppositions, such antinomies, as require a \emph{mediation in the concept}; for, as it is said in the Wisdom of Sirach (33:17): ``The works of the Most High are always two, one over against the other''; and the speculative rests precisely upon grasping the oppositions in the unity of the idea. A speculative beholding must always be presupposed, if the presentation is not to lose itself in outward understanding, or confine itself merely to apprehending the Christian conceptions in their practical, merely applicable significance.'' --- Therefore, too, shall all those have no thanks, who (p. 85) ``distinguish between theology and philosophy, as between clean and unclean foods, who say of philosophy: Taste
% --- p. 71 ---

not, touch not! not considering, that in their own theology, too, which is, after all, likewise in many respects a work of man, there might --- be much that is unclean, that can only be purified by philosophy''; not considering, that it is the same Logos, that works in the kingdom of nature and in the kingdom of grace'' etc.

Since it is, then, profoundly said, the same Logos, that works in the kingdom of nature and in the kingdom of grace, and since dogmatics, moreover, prosaically said, can quite well make use of philosophy's ``logical and ontological investigations,'' as a ``propaedeutic foundation,'' ``the Christian Dogmatics'' will, then, in view of this, with Christian charity, take under its wing the poor, forsaken, and wayward speculation, well knowing, that it is by no means% <-- print "ev" (r/v glyph); corrected
{} Christianity, but speculation, that is hereby done a great service, well knowing, ``that it is by no means religion, that is to borrow its weight and its significance from speculative thinking, but that it is speculative thinking, that stands in need of religion, of God's revelation, as of its principle.''

Before we, however, give ourselves over to joy, and congratulate ``speculative thinking'' on its dogmatic exaltation, we ought, nonetheless, for safety's sake, to look somewhat more closely, whether the conditions are, indeed, tolerably favorable, whether ``speculative thinking'' can really, in fact, be served by the terms offered, for --- since I am, indeed, now also to say something with definiteness --- many and remarkable are the passages, that, in this ``Christian Dogmatics,'' speak to the abasement of speculation.
% --- p. 72 ---
\emph{``The dogmatician (p. 4) is only the organ of his science, in that he is, at the same time, the organ of his Church, which does not hold of the mere philosopher, who will serve pure science alone.''}

Already this declaration awakens, in a high degree, the mistrust of speculation, that with ``coherent thinking'' things do not, after all, quite hang together properly. What does the author now mean here by ``pure science''? According to Hegelian usage, one would most likely think of logic. But the mere philosopher does not, indeed, confine his speculation to pure logic alone; it is, on the contrary, indeed, presupposed (p. 81, § 35), ``that there is such a thing as a Christian philosophy,'' it is presupposed, ``that this philosophy is bound to the same fundamental conditions of cognition as theology, and must, therefore, proceed from \textit{credo ut intelligam}.'' By the expression, ``pure science,'' one must, then, here reasonably think of the science, that has the same content as dogmatics. But if the expression is to be understood in this way, speculation must, indeed, both on behalf of science in general, and on its own behalf in particular, guard itself well against so ambiguous an adoption. The dogmatician announces his dogmatics as science. However this science may, for the rest, be determined, philosophical or not philosophical, the dogmatician, as the organ of his science, can, nonetheless, possibly have no other task than that of satisfying the demand of his science; and yet it is, nonetheless
% --- p. 73 ---
emphasized as a significant nota bene, that the dogmatician is only the organ of his science, in that he is, at the same time, the organ of his Church. How is this? --- speculation asks --- does the dogmatician, then, now speak with tongues? Is being the organ of one's science the same as being the organ of one's Church, what, then, does that nota bene mean? But if being the organ of one's science is one thing, and being the organ of one's Church another: ought one not, then, first to be informed a little more closely, how two tasks so dissimilar let themselves be united into one task? If the dogmatician, in being the organ of his science, is, at the same time, to be something other than the organ of his science, then he must, indeed, make his science into something other than science, that is to say, treat his science in a supra-scientific, which no doubt means: an unscientific way. Would it not, then, be better, if the dogmatician, as the organ of his science, like the mere philosopher, first of all set himself to serve his pure science alone; should it then, alongside this, turn out, that the Church required yet another service of him: ``the mere philosopher'' would have nothing to object. For that a sincere dogmatician, ``in another way too'' than as dogmatician, seeks to benefit his Church, speculation can, then, least of all take amiss, since ecclesiastical practice, indeed, at bottom, has nothing whatsoever to do with pure knowledge.

As regards speculation's own interest, the objection is, in short, this: to apply speculation to dogmatic-ecclesiastical use is to use it for some-
% --- p. 74 ---
thing else, to displace its purpose outside itself, to speculate relatively and to a certain degree, that is: to speculate outside speculation. For speculation thinks the Absolute absolutely, and has, as such, its purpose in itself alone. Speculation is an ideal potency, that must stand or fall by its own strength, i.e. by its own relation to the Absolute. One can declare oneself free of speculation, one can pass it by, overlook it, despise it; but one cannot use it for something else, cannot apply it relatively, without violating the Absolute, that is seen in its medium and will take shape in its dialectic. How, then, does the dogmatician manage to become the organ of speculation, while at the same time wishing to be the organ of his Church?

\emph{``For dogmatics,'' --- it says (p. 5) --- ``Christianity's absolute truth is given in advance and independently of all speculation. That \textit{δός} πᾷ στῶ, which the seeking philosophy has so often uttered, is originally answered for dogmatics; and the dogmatician does not make Christianity dependent on his own inquiry, but seeks, through thinking, only a deeper appropriation of the truth, that is, to him, the surest of all things, and to which he has come by a wholly different way than that of speculation.''}

Against the first half of this quotation, speculative thinking has nothing to object; on the contrary, it
% --- p. 75 ---
must even acknowledge, that dogmatics, on these terms, is the happiest, best-endowed, and wealthiest of all the sciences, since it has, from the cradle, everything, that the other sciences must laboriously work for. See, it has, indeed, the standpoint, or, more precisely, it has that \textit{δός} πᾷ στῶ, that the seeking philosophy has so often lacked; it has the truth, or, more precisely, it has Christianity's absolute truth; and all this is given to dogmatics ``in advance and independently of all speculation.'' But, when dogmatics really has these fair gifts in advance and independently of all speculation, one asks, with good reason, what more it wants, then, why it, then, moreover, still wishes to condescend to desire, what belongs to speculation? It is by no means because speculation is envious, that one asks thus: speculation gladly grants dogmatics all possible honor and prosperity, would also gladly congratulate the dogmatician both on his πᾷ στῶ and on the absolute truth; but one asks thus, because one does not, in any way, comprehend, what moves the so richly endowed% <-- print "boterede" (d/b glyph); corrected
{} dogmatics to enter into partnership with the poverty of speculation. Good heavens, when dogmatics already, beforehand, has both πᾷ στῶ and Christianity's absolute truth, then it might, indeed, gladly be content; and the Christian dogmatician ought, surely, not to hold it against speculation, that it, here, in relation to so superfluous a wealth, considers its own poverty, and holds itself back a little, as though it were afraid of being left over.
% --- p. 76 ---
The significance attached to speculation in the quotation's latter half is exceedingly mysterious, indeed so mysterious, that I, for my part, should not be disinclined to hide my embarrassment behind the well-enough-known phrase: ``I acknowledge the mystery,'' if only, by such a phrase, one could get well off from a half-completed inquiry. ``The dogmatician does not make Christianity dependent on his own inquiry, but seeks, through thinking, only a deeper appropriation,'' etc. What kind of \emph{thinking} is this, now? So far as I can judge, there is here but one way out to rescue any meaning in this dark passage, and that is, to explain the dogmatician's \emph{thinking} not as a speculative, but as an existential thinking. Is this permitted? There are, indeed, to be found here and there in the Introduction, a good many passages, that speak for such a conception. It says (p. 9, § 5), that the religious relation to God must, more precisely, be determined as ``the holy, personal relation to God, that has its general expression in the relation of conscience,'' further (§ 5), that conscience is ``man's original co-knowledge with God concerning his personal relation of existence to him''; further, that ``as I know myself in my conscience, so do I live and exist''; further: that ``all certainty concerning divine and human things is only a religious certainty, when it is a certainty of conscience''; --- when all these propositions are duly weighed and thought through, we do, indeed, seem undeniably to arrive at the result, that the dogmatic thinking, of which mention is here made, must
% --- p. 77 ---
be the existential, and the appropriation likewise existential. Might we now but assume this, then the whole matter would be settled at one stroke; dogmatics would then, in a single vigorous turn, have taken leave of speculation, and so, full stop. But it will not do. How should ``the Christian Dogmatics'' take leave of speculation (p. 5); when the dogmatician himself, as we have already seen it in so many other passages, both vigorous and beautiful (pp. 78, 79, 85, and elsewhere), with definiteness demands speculation, and believes in it, and bears witness for it. For further confirmation, we will once more call to mind the testimony borne in the preface (p. II): ``That a coherent theological thinking, in itself, indeed theological speculation, is both possible and necessary, is a conviction, I hope to be able to hold fast, even should it be demonstrated, that my own work had failed.'' --- See, that I call \emph{believing} --- in speculation.

``\ldots{} and the dogmatician does not make Christianity dependent on his own inquiry, but seeks, through thinking (that is to say, and let it be written down: through speculative thinking) only a deeper appropriation of the truth, that is, to him, the surest of all things, and to which he has come by a wholly different way than that of speculation.'' By which other way? --- ``The way of faith.'' Therefore: the dogmatician has, by the way of faith, already arrived at this, that Christianity's absolute truth is, to him, the surest of all things! --- If the dogmatician were now a plain single individual
% --- p. 78 ---
in altogether the same way, as every other believer in the congregation is a single individual, then, indeed, ``the Christian Dogmatics'' would, presumably, first cast upon him a grave look, and ask him, in its preface, whether he, the single individual, had not, after all, become ``all too soon'' ``absolute in faith.'' But since ``the dogmatician'' is not a plain single individual, like other ordinary single individuals, but stands in the Church, and is only the organ of his science, in that he is, at the same time, the organ of his Church; since he, in his capacity as dogmatician, has fused, or at least as good as fused, with ``that great single individual''; well, then, he shall also, so far as the principle is concerned, be permitted to begin as absolute in faith. As absolute in faith, the dogmatician has, then, come to Christianity's absolute truth by a wholly different way than that of speculation. Yes, that is correct, that is altogether clear and correct! but we must now pause a moment, for now the mystery comes, and I have, indeed, admitted, that I ``acknowledge the mystery.'' If Christianity's truth is, to the dogmatician, the surest of all things, and if he has really come to the cognition of this truth by a wholly different way than that of speculation's; but what in all the world does the dogmatician want, then, with speculation? --- ``He seeks only, through thinking, a deeper appropriation.'' --- Through thinking? Is there, then, no thinking at all in faith, or is all the thinking of faith, at bottom, speculative thinking? Here it does, indeed, become necessary to answer with ``a dry yes or no''; for even if one wished to answer with a fluid approximation, that the thinking of faith already from
% --- p. 79 ---

its first germ was half and half a speculative thinking, I am, nonetheless, afraid, that one would get something quite other out of it, than what the dogmatician wishes. How, in general, am I to assure myself of a truth, except through a thinking? The idea of truth is, indeed, after all, as is well known, the idea of thought, and a thoughtless assurance of the idea of truth is, consequently, altogether untrue. We must by no means imagine, that Dr. M. should be ignorant of this. It is true, he says in the preface, ``that faith is the one thing needful for us all,'' and that ``those, then, who feel no need to reflect upon their faith, may well be within their rights, when they, for their own part, regard coherent cognition as unnecessary.'' But this is, nonetheless, surely to be understood \textit{cum grano salis}; it is, so far as I can judge, pure irony at the expense of the poor single individuals, who now also wish to have their say, though they, indeed, know as little about faith as about knowledge. That what is said in the preface really is irony, I conclude, moreover, from yet another reason. Dr. M. says, first, that faith is the one thing needful for us all, and lets it now appear, as though he meant, that thinking was needful only for those, who feel a particular need for coherent cognition, that is to say, for the science of faith and speculation, and thus needful only as an exception to the general rule. But next he censures it, both as ``a new and an old misdirection,'' that some single individuals should hit upon wishing to separate faith from coherent cognition, and now lets it appear,
% --- p. 80 ---
as though he meant the opposite, that not only faith, but also the cognition of faith, is ``the one thing needful for us all,'' and that, accordingly, those single individuals, who either cannot or will not enter into the science of faith and speculation, are by no means within their rights, but can only, at a pinch, be tolerated in the congregation as an exception to the general rule. So open a contradiction, that, in one and the same breath, makes rule into exception and exception into rule, would, indeed, be altogether intolerable, were the whole not taken as a feigned parade, a sarcastic half-jest, an irony. If we now abstract from the irony, then it stands, indeed, immovably firm, that Dr. Martensen, with definiteness, acknowledges the indissoluble unity of faith and thinking. In this conviction, we then presuppose as settled, that the truth of faith can only be appropriated through a thinking; and thus the question returns: is the thinking of faith different from that of speculation, or are they, at bottom, one? If the dogmatician grants the latter, then we do, indeed, concede, that he is right, when he seeks the completion of the cognition of faith in speculation; but, we must then ask, what entitles him to the claim, that he himself has come to the highest certainty of Christianity's absolute truth by a wholly different way than that of speculation? If, on the other hand, he declares himself for the first assumption, namely that the cognition of faith is, in principle and essence, altogether different from speculative cognition, then we gladly concede to him, that he must, in advance, have come
% --- p. 81 ---
to the absolute truth by a wholly different way than that of speculation; but we must then ask, where he, as dogmatician, then wishes to go with speculation? --- ``He wishes to see the glory of God.'' --- Where? --- ``In speculation! For let Melanchthon say what he will against idle speculation: the cognition of God's mysteries is, nonetheless, in itself, a good, and the cognition of God's glory, in itself, edifying!'' --- Yes, assuredly edifying! But the question is, whether this edifying cognition is to be found in speculation. What does Christ say to Martha? If thou couldst believe, thou shouldst see the glory of God (John 11:40). But Martha's vision was, nonetheless, assuredly no speculative vision, which can be concluded not merely from the fact, that she was a woman, but best of all from the fact, that she, the sorrowing one, saw through tears, when she saw the glory of God; and when the believer sees the glory of God through tears, the vision is, indeed, edifying, but not speculative. --- ``Yes, but God forbid, there is, indeed, no one, either, who would deny that; but, because Martha saw the glory of God without speculation, a Christian dogmatician may, therefore, nonetheless, well feel a need to see the glory of God in speculation, too; for, as it is said in the Wisdom of Sirach (33:17): `The works of the Most High are always two, one over against the other,' and the speculative rests precisely upon grasping the oppositions in the unity of the idea.'' --- But even granting the wise Sirach be right, when he says, that the works of the Most High are always two, one over against
% --- p. 82 ---
the other; is it, therefore, said, that the dogmatician can be right, when he wishes to mediate the works of the Most High, or is it not, rather, precisely the case, that no one can mediate the works of the Most High but the Most High himself? --- ``It is a pure misunderstanding, if anyone should interpret the author's words as though the meaning were this, that the dogmatician now really presumed to see through God's mysteries and mediate all God's works. On the contrary, the dogmatician acknowledges, indeed, the mystery! The dogmatician will, therefore, also only speculate to a certain degree, and mediate to a certain degree; but nonetheless --- let it be said with definiteness --- to a certain degree. We must, indeed (p. 79, § 33), always hold fast, that comprehending is the \emph{piecemeal} in our cognition, while the whole lies in the fullness of intuition, that is not exhausted by any development of the concept. But just as those have always been put to shame, who have given themselves out as having comprehended everything, so, too, have those been no less put to shame, who have thought to fix, once and for all, a \textit{non plus ultra}, whereby the matter should have its final settlement. For, afterward, it has always shown itself, that there was, nonetheless, a `\textit{plus ultra}' after all, and the supposedly fixed boundaries revealed their mobile nature by shifting themselves.'' ---

Yes, by shifting themselves! The boundary shifts, the boundary is fluid. We have, indeed, already heard, that here, on account of the boundary's being fluid, one no longer dares speak absolutely of the Absolute. But let us, then, divide the in-
% --- p. 83 ---
quiry, and say a word or two, first about \textit{plus ultra}, and then about mediation. That speculation finds itself in a steady advance can easily be shown from the history of philosophy. Thus Fichte, for instance, is \textit{plus ultra} in relation to Kant, Schelling in relation to Fichte, Hegel in relation to Schelling; but who, then, is \textit{plus ultra} in relation to Hegel? A good many aspirants have, indeed, presented themselves of late, but none of them has, so far as is known to me, yet won the prize. Hegel stands, then, still, in pure speculation, as the unvanquished master, as \textit{non plus ultra}. How, now, does ``the Christian Dogmatics'' judge of Hegel's mediation? We read, on p. 84: ``An uncritical intermixture of dogmatics and philosophy we find, thus, in many ways, among the Alexandrian theologians, where the categories of Platonism often stepped in for those of Christianity. The same thing was repeated in the Middle Ages under Aristotelian influences; and it is still fresh in memory, how the more recent Aristotle, how Hegel, occasioned the same thing. In the false mediation, in the sham reconciliation between faith and knowledge, one is reminded of Augustine, who, in his \textit{retractationes}, confesses, that in his Platonic period he read Plato out of the Gospel, and thus thought to have found the reconciliation of religion and philosophy.'' --- Let us take note of these words, for they are well worth our laying them to heart. If Hegel is still, in speculation, \textit{non plus ultra}, and if it is, nonetheless, this very philosopher, who, by his
% --- p. 84 ---
``false mediation,'' by his ``sham reconciliation between faith and knowledge,'' has brought about the more recent age's ``uncritical intermixture of dogmatics and philosophy,'' then, ye immortal gods, are we not, then, as regards the speculative comprehending of Christianity, still exactly where we were! And yet, perhaps I understand it. Hegel was put to shame, because he could only \emph{think} speculation: since that time, one has learned something definite from ``the signs of the times,'' and found \textit{plus ultra} in \emph{believing} in speculation; Hegel was put to shame, because he still knew only a speculation indifferent to the edifying; since that time, one has learned to add something, and found \textit{plus ultra} in \emph{edifying} speculation; Hegel was put to shame, because he mediated the relative as absolute; since that time, one has learned to strike something off the demand of mediation, and found \textit{plus ultra} by mediating the Absolute relatively. Well: this I understand; but, once one has, indeed, given up mediating the Absolute in the absolute concept, is it, then, also certain, and in every way settled, that ``the works of the Most High'' absolutely require a mediation in the relative concept? ---

``Yes, that is altogether certain. A speculative \emph{beholding} must always be presupposed, if the presentation is not to lose itself in outward understanding, or confine itself merely to apprehending the Christian conceptions in their practical, merely applicable significance.'' --- Therefore: when the Christian conceptions are merely appropriated in faith, can they, then, do nothing but lose themselves in an outward understanding?
% --- p. 85 ---
--- ``But what a misunderstanding! Who could own to so unreasonable a notion? Is it, then, not expressly said, and said with definiteness, that the dogmatic cognition is itself a cognition in faith and out of faith? Is it not said, that the religious relation to God must, more precisely, determine itself as the holy, personal relation to God, that has its general expression in the relation of conscience? Is it not said, and said with definiteness, that conscience is man's original co-knowledge with God concerning his personal relation of existence to him, that, as I know myself in my conscience, so do I live and exist? Is all this, and much more, not said in the Christian Dogmatics, and said with definiteness?'' --- If I have, indeed, on so important a point, really misunderstood ``the Christian Dogmatics,'' I beg pardon most sincerely, and now wish only to know, why the dogmatician cannot confine himself to apprehending the Christian conceptions in their ``practical and merely applicable significance''? --- ``This question is, then, easily answered: the dogmatician seeks a deeper appropriation of the truth, that is, to him, the surest of all things; this deeper appropriation he cannot attain without a speculative beholding, and therefore he makes use of speculation for a scientific presentation and grounding of the Christian articles of faith.'' --- Hear a word: every grounding of truth must, surely, indeed, be designed to heighten certainty and strengthen conviction; but an appropriation of the truth, that is, to one, the surest of all things, can, nonetheless, surely, never become
% --- p. 86 ---
deeper, in relation to grounding and certainty, than it already is; unless we should someday experience the curiosity, that what is, to the dogmatician, the surest of all things, should, in the end, become to him even surer than the surest of all things. --- ``A fruitless taunt! there is here, indeed, no talk at all of the dogmatician's practical certainty in faith.''\ldots{} --- No, that is, assuredly, as great as it can be.\ldots{} --- ``But the talk is of a heightening of the theoretical and speculative insight into the Christian mysteries.'' --- If it now really becomes a serious matter with the speculative insight into the Christian mysteries, then the dogmatician will, indeed, also be obliged to follow speculation's instruction, and hear the demands, that objective thinking lays upon the one who speculates.

In the first place, then, independently of all practical considerations, a direct relation must be maintained between theoretical certainty and speculative insight: the more insight, the more certainty; the less insight, the less certainty; absolute insight, absolute certainty. Only in this way can the speculative problem arise. Let us illustrate this with a couple of examples. The Eleatics knew, indeed, surely, just as certainly as other men, that in the sensible world there is a becoming and a movement. They had, concerning this, an immediate and practical certainty; but nonetheless, they denied to becoming and movement all essential truth. Whence, now, did that come? It came precisely from this, that they could not think these concepts speculatively. The conception of a becoming was
% --- p. 87 ---

in contradiction with their pure insight, and, in virtue of this insight, they denied the reality of the conception. The subjective idealists know it just as fully as other reasonable men, that in an immediate, practical sense there is an outward, sensible world with things and objects; nonetheless, they treat, in their speculative thinking, this outward world as a world of appearance, a barrier, a not-I, that has significance only for the I. Whence, now, does that come? It comes from this, that the conception of the outer world's reality is in conflict with their pure insight. Why does Schelling go back to the indifference of the ground? Why does Hegel begin with being and nothing? Why does Sibbern define philosophy as an all-encompassing discussion and debate of everything against everything? Naturally, because speculation, as speculation, does not concern itself at all with practical certainty, but works solely and alone to attain theoretical certainty in the pure insight. If, then, speculation is to be carried through in dogmatics, the dogmatician must begin from the ground up, and transform all Christianity's practical conceptions into speculative problems. But so bold an operation the dogmatician surely dares not risk; for what would then become of ``the Christian Dogmatics' cognition in faith and out of faith''?

The next thing the one who speculates has to observe is the dialectical relation between presupposition and abstraction; for on this, indeed, precisely depends, what one, in
% --- p. 88 ---
more recent times, has called the method.\footnote{Cf. my work, \textit{The Propaedeutic Logic}, pp. 102--112.} The question, whether speculation, and \textit{in specie} the speculation of Christianity, should begin with or without a presupposition, has, in more recent times, become so worn threadbare, indeed so hackneyed, that one can scarcely feel anything but dread at touching it, especially when one, as is the case in this context, has only a couple of general propositions at one's disposal. So much can, nonetheless, be said in all brevity, that the vulgar question of presupposition or no-presupposition is empty and meaningless. That the one who speculates must be furnished with knowledge and presuppositions corresponding to the subject he wishes to treat, goes, indeed, without saying. Or how, indeed, should anyone be able to speculate on Christianity without Christian knowledge and Christian presuppositions? But the point is, that the one who speculates must begin by transforming all his presuppositions into speculative problems, all dogmas into philosophemes; without such a radical transformation, nothing whatsoever comes of speculation. In that the presuppositions are now transformed into problems, one does not merely abstract from their purely practical validity, but the abstraction seeks back to the dialectical point, where the one who speculates has, in conception, everything, and yet, in objective knowledge, still nothing, i.e. nothing other than speculation's own first breath and methodical beginning. In this
% --- p. 89 ---
sense, the speculative philosophers, both Schelling and Hegel, have had Christian presuppositions, and if their philosophy of religion has, nonetheless, not yielded the expected results, the reason for this must not be sought, in the first instance, in a lack of presuppositions, but far rather in a deficiency in speculation itself. That this is so can, moreover, be proved from the example of the theologians who followed in the tracks of those philosophers. In speculative cognition, the dogmaticians Daub and Marheineke have essentially not carried it further than the philosophers Schelling and Hegel. But now, surely, no one will maintain, that Daub and Marheineke lacked either theological knowledge or Christian presuppositions. When, now, nonetheless, a Christian dogmatician, who believes --- in speculation, and bears witness for ``coherent cognition,'' confuses% <-- print "forverler" (r/x glyph); corrected
{} practical certainty with theoretical, speculative abstraction with lack of presupposition, overlooks Hegel, pushes Marheineke into the background, calls the speculative relation to God ``a relation at second hand'' (p. 8, § 4), speaks pathetically of ``faith's eternal power of youth, its power to unfold, out of its own depths, a wealth of the treasures of knowledge and cognition''; while he, nonetheless, at the same time, feels a need ``through speculative thinking to seek a deeper appropriation of the truth, that is, to him, the surest of all things'': then one hardly knows what to believe, or, rather, then one, unfortunately, knows only too well what to think of the use, ``the Chris-
% --- p. 90 ---
tian Dogmatics'' here intends to make of speculation.

If faith has, in itself, an eternal power of youth, a power to unfold, out of its own depths, a wealth of the treasures of knowledge and cognition, while the speculative relation to God nonetheless always is and remains ``a relation at second hand'': would it not, then, be far better, if the dogmatician gave up the speculative beholding of God's glory, and made do, instead, with the thinking of faith, in order, thereby, to render ``the Christian conceptions in their practical, merely applicable significance''?

\begin{center}
\emph{c) ``The Christian Dogmatics'' will appropriate faith, but without holding fast to faith's problem.}
\end{center}

It cannot be held against the author of ``the Christian Dogmatics,'' that he in any way seeks either to undermine, or willfully to disparage, faith. He is no dryly acute rationalist, who, by all manner of reasonings, explains away Christianity, until he has, at length, explained away the Christian element in its principal doctrines. Nor, either, is he an enthusiastic thinker, who, working only for the concept, dissolves the dogma and volatilizes the mystery into pure knowledge. On the contrary! The dogmatician is, indeed, ``only the organ of his science, in that he is, at the same time, the organ of his Church.'' As the organ of his Church, he does, indeed, to be sure, proceed from the presupposition, that speculation, at least a certain speculation, must necessarily develop out of the doctrine of faith, roughly in
% --- p. 91 ---
the same way, that dogmatics itself develops out of the catechism; but this presupposition is altogether naive, it comes as though of itself, and is assumed \textit{de bonne foi}, without any doubts or skeptical difficulties. This naivety is, now, so far a gladdening sign, as it bears witness, that the dogmatician does not speculate, in any way, in any unecclesiastical or unchristian sense, but believes in speculation solely because he believes, that it is for ``the glorification of faith, is \textit{in gloriam fidei, in gloriam Dei}.'' Should it then, someday, show itself, that dogmatic speculation, according to its nature, is neither for God's honor nor for the glory of faith: then there is, indeed, no doubt either, that the Christian dogmatician will at once abandon speculation, and hold fast to faith alone. ``For no price,'' --- says Dr. Martensen in his preface --- ``would I be at variance with the believer, and willingly shall I give up any of my propositions, when it is shown me, that they really entail such a variance.'' See, that is an open confession, a benevolent concession, and a good profession, upon which much can, in time, be built.

Meanwhile, the author's naive faith in speculation's dogmatic usefulness has, nonetheless, all the same, a questionable shadow side. For, had the dogmatician not, with quite so naive a certainty, confounded the existential and the speculative thinking with one another: ``the Christian Dogmatics'' would, indeed, most assuredly, have contained a good many inquiries, that the thinking reader must now, unfortunately, sorely miss.
% --- p. 92 ---
There are, principally, two cardinal propositions, to which ``the Dogmatics'' refers us, when the talk is of the necessity of a speculative cognition and an objective knowledge. The one of these propositions is the Anselmian: \textit{credo, ut intelligam}. ``By its \textit{credo, ut intelligam}'' --- it says (p. 73, § 33) --- the Christian dogmatics distinguishes itself from the cognition, that proceeds from the proposition \textit{de omnibus dubitandum est}, insofar as it is hereby required, that thinking shall loose itself from all presuppositions, and go out on voyages of discovery, in order to find the truth, whatever it may be. Christian cognition receives its impulse not from doubt, but from faith. There can, on the other hand, indeed, be talk of a skepticism in dogmatics, insofar as such a skepticism is only the expression of the critical and dialectical drive contained within faith.''

Had the dogmatician now but made earnest of this ``skepticism,'' and let dogmatics develop ``the critical and dialectical drive contained within faith'': then he would, assuredly, also have discovered the principled difference, that exists between the cognition of faith and speculation. After what has, concerning this matter, already been set forth in the Review of the Pseudonym, I consider it superfluous to treat the matter afresh; only in order to point back to what has already been said, and thereby, at the same time, to call to mind the ingenuity, with which the indefatigable dialectician understands how to nourish the inquiry, and keep the problem
% --- p. 93 ---
alive, do I here cite a couple of passages verbatim from \emph{Johannes Climacus}.

``That the speculative view is objective is not denied; on the contrary, merely in order further to demonstrate it, I will here again repeat the attempt, of setting the subjectivity, infinitely concerned in passion for its eternal blessedness, in relation to it, whereby it will then, precisely, become manifest, that the speculative view is objective, in this, that the subject becomes comic. He is not comic, because he is infinitely interested (on the contrary, everyone is, indeed, comic, who, without being infinitely interested in passion, will nonetheless wish to make people believe, that he is interested in his eternal blessedness); no, the comic lies in the disproportion of the objective.''

``If the speculator is at the same time a believer (which is also said), then he must, indeed, long since have realized, that speculation can never have for him the significance that faith has. Precisely as a believer, he is, indeed, infinitely interested in his eternal blessedness, and is, in faith, assured of it (N.B. as one, believing, can be assured of it, i.e. not once for all, but daily acquiring, with the infinite, personal, passionate interest, the certain spirit of faith); and he builds, then, no eternal blessedness upon his speculation, he deals, rather, suspiciously with speculation, lest it should trick him out of the certainty of faith (which, in every moment, has within it the infinite dialectic of uncertainty) into the indifferent objective knowledge. Simply, dialectically understood, the matter stands
% --- p. 94 ---
thus. If he therefore says, that he builds his eternal blessedness upon speculation, then he contradicts himself comically, for speculation, in its objectivity, is precisely altogether indifferent to his and my and your eternal blessedness; while the eternal blessedness lies precisely in the self-diminishing self-feeling of subjectivity, acquired by the utmost exertion. He lies, too, with respect to his giving himself out as a believer.''

``Or the speculator is not a believer. The speculator is, naturally, not comic, for he does not ask at all about his eternal blessedness; the comic arises only, when the subjectivity, infinitely interested in passion, will set its blessedness in relation to speculation. The speculator, on the other hand, does not set forth the problem, we speak of; for, as a speculator, he becomes, precisely, too objective to trouble himself about his eternal blessedness. Here, though, one word, so that it may become clear, should anyone misunderstand many of my utterances, that it is he who will misunderstand me, while I am without fault. Honor be to speculation, praised be everyone who truly occupies himself with it! To deny the worth of speculation (even though one might wish, that the money-changers in the forecourt were driven out as profane) would, in my eyes, be to prostitute oneself, and especially foolish of one, whose greater part of life, as poor opportunity allowed, has been dedicated to its service, especially foolish of one who admires the Greeks. For he must, indeed, surely know, that Aristotle, when he discusses what blessedness is, places the highest blessedness in
% --- p. 95 ---

thinking, reminding us of the blessed pastime of the eternal gods in thought. He must, moreover, have both a conception of and a reverence for the undaunted enthusiasm of the man of science, his perseverance in the service of the idea. But for the one who speculates, the question of his personal eternal blessedness cannot come forward at all, precisely because his task consists in coming, more and more, away from himself, and becoming speculation's beholding power.''

``Understood in this way (and, in a certain sense, it would not even need to be pointed out, that the infinite interest in one's eternal blessedness is something higher, for the main point is, that it is the thing asked about), the comic will readily show itself in the contradiction. The subject is infinitely interested in passion for his eternal blessedness, now he is to be helped by speculation, that is to say: by himself speculating. But, in order to speculate, he must, precisely, go the opposite way, give up and lose himself in objectivity, disappear to himself. The disparity will altogether prevent him from beginning, and will, comically, judge any assurance of having gained anything by this way. The contradiction between the subject, infinitely interested in passion, and speculation, when it is to help him, I will permit myself to illustrate by a sensible image. When one wishes to saw, the point is, not to press too hard upon the saw; the lighter the sawyer makes his hand, the better the sawing goes. Should one press with all his might
% --- p. 96 ---
upon the saw, he does not get to saw at all. So it holds, too, here, for the one who speculates, that he must make himself objectively light, but the one who, in passion, is infinitely interested in his eternal blessedness, makes himself as subjectively heavy as possible. Precisely thereby he makes it impossible for himself to come to speculate.'' --- If this now stands so; then the dogmatician's ``\textit{credo, ut intelligam}'' can, indeed, in no way justify speculation.

``The Dogmatics''' second main proposition aims at forestalling (p. 77) ``the misdirection, that, out of concern for `the believer' and his state of soul, becomes indifferent to faith (\textit{fides, quæ creditur}), that, out of zeal for the edifying, forgets the content, with which one is to build, forgets the foundation, upon which one is to build''; --- this misdirection the dogmatician now believes he can forestall, by emphasizing both the objective (\textit{fides, quæ creditur}) and the subjective (\textit{fides, qua creditur}: \textit{res ipsa} and \textit{usus}). Catholicism inclines too much to the objective (\textit{fides, quæ creditur}); one-sided Protestantism inclines too much to the subjective (\textit{fides, qua creditur}): ``the Christian Dogmatics'' holds the balance between the objective and the subjective, and precisely therein consists the truth. ``What the Reformation wanted,'' --- it says (p. 40, § 21) --- ``was neither exclusively the objective nor the subjective; it was the free union of the objective and the subjective, of the content of faith and the inwardness of faith, of divine revelation and religious self-consciousness.''
% --- p. 97 ---
All this may, now, be quite sensible, what is said thus, in general, about ``\textit{fides, quæ creditur}, and \textit{fides, qua creditur}'', as well as about ``the free union of the objective and the subjective,'' if one could only get to know, wherein this union is, then, really to consist, and how far it lets itself, in reality, be realized; but how, indeed, should this be developed in a dogmatics, where the problem of inwardness does not occur at all. Let us, for once, consider the doctrine of the sacraments. This doctrine does, indeed, assuredly presuppose a free union of the subjective and the objective: of what nature, now, is this union?

Of baptism it says: ``He that believeth and is baptized shall be saved; but he that believeth not shall be damned'' (Mark 16:16). That this proposition is, in and of itself, an objective doctrinal proposition is proved by this, that the Church exists, and that the existing Church confesses it. The assurance of this is, then, to that extent, objective, as it rests immediately upon a facticity. The non-believer can have this assurance just as fully as the believer. The subjective is, then, here, the indifferent: the matter is decided solely by the factual objectivity. But it is another question, whether one can, by a purely objective consideration, assure oneself of the objective truth of the doctrinal proposition, whether one can, then, in this sense, first determine and settle the content of faith (\textit{res ipsa}; \textit{fides, quæ creditur}), in order thereafter to arrange and regulate the subjective
% --- p. 98 ---
faith (\textit{usus}; \textit{fides, qua creditur}). This question must be answered in the negative. The proposition, ``He that believeth and is baptized shall be saved,'' is only objectively true for the believer; and, likewise, the proposition, ``He that believeth not shall be damned,'' is only objectively true for him who believes. Were it possible, by any proof whatsoever, whether historical or speculative, to establish in advance the objective truth of the doctrinal proposition: then this truth would, indeed, have to be equally evident to the non-believer as to the believer; but if a non-believer could, in this way, be convinced of the truth of revelation, he would, indeed, by following the proof, necessarily have to become either a believer or a devil. One must not now seek evasions, and say, that the objective grounding, naturally, only holds to a certain degree, and always points on to the subjective; for, when one says, that the objective grounding holds to a certain degree, one incurs the misunderstanding, that it, nonetheless, always accomplishes something for faith. No, one must go quite plainly to confession, and admit outright, that what the objective grounding here accomplishes is altogether nothing. But if the objective grounding can accomplish nothing, nothing whatsoever: what, then, follows from this? Well, then it follows, indeed, from this, that the assurance of the objective truth (\textit{fides, quæ creditur}) must rest exclusively upon the subjective grounding, i.e. upon the existential appropriation in faith (\textit{fides, qua creditur}). In the same degree that I, thinkingly, understand, and, existingly, ex-
% --- p. 99 ---
press, what it is \emph{to believe}, in the same degree does, too, my essential conviction and inward assurance of the truth rise, in the revealed word: ``He that believeth and is baptized shall be saved''; but everyone, who truly believes, that he shall be saved by faith, must, indeed, necessarily also believe, that he would, by not believing, become unblessed, that is, damned. --- The doctrinal proposition's latter clause stands, then, in the most inward connection with the first; but, by the objective consideration, and the objective proof, and the objective grounding, the whole is thrown into confusion.

Of the Supper, it is said by the Apostle Paul (1 Cor. 11:28--29): ``Let a man examine himself, and so let him eat of that bread, and drink of that cup; for he that eateth and drinketh unworthily, eateth and drinketh damnation to himself, not discerning the Lord's body.'' --- Of the objective truth in these words, no dogmatician can assure himself objectively; he can, on the other hand, by objective proof and objective grounding --- yes, this shall not be denied --- displace the point of view, and throw the believing appropriation into confusion. The understanding of Paul's words is not partly objective and partly subjective, but solely subjective, absolutely subjective. He who will not rightly examine himself, before he eats of the bread and drinks of the cup, cannot, by any objective proof, assure himself, that he, as an unworthy one, will indeed eat and drink damnation to himself; for his unworthiness lies, indeed, essentially, in his unbelief, and his
% --- p. 100 ---
unbelief consists, indeed, precisely in this, that he does not properly ``discern the Lord's body.'' Only the believer, who examines himself, in order to eat and drink worthily, understands, in his conscience, how terrible it is, to eat and drink unworthily, that is, without discerning the Lord's body.

``But,'' --- one says --- ``the significance of the subjective is, indeed, also fully acknowledged in the Christian Dogmatics; or how, indeed, should dogmatics be able to impress subjectivity and the inwardness of faith more emphatically, than when it is said, that the dogmatic cognition is a cognition in faith and out of faith; that the religious relation to God must, more precisely, be determined as the holy, personal relation to God, that has its general expression in the relation of conscience; that conscience is man's original co-knowledge with God concerning his personal relation of existence to him; that, as I know myself in my conscience, so do I live and exist: surely no one, then, will demand a greater acknowledgment of the principle of subjectivity?'' --- No, it would be a shame to say, that the acknowledgment, in these propositions, is not great enough; but, when the dogmatician has, then, thus, acknowledged the principle of subjectivity: why, then, does he not abide by this acknowledgment, and prove his faith by his works? why, then, does he go over into an acknowledgment of the altogether opposite, in altogether opposite propositions? --- ``The works of the Most High are always two, one over against the other.'' --- Yes, the works of the Most
% --- p. 101 ---
High! --- When the dogmatician gets onto the subjective track, he acknowledges and impresses subjectivity in the best way; when he then, again, falls back into the speculative train of thought, he maintains objectivity, objective knowledge, ``theological speculation,'' in the most emphatic way; and yet these two directions are not merely altogether opposed to one another, but also altogether contradictory.

But, because the believing appropriation is absolutely subjective, it by no means follows from this, that it should exclude the objective. It is so far from the case, that the subjective faith (\textit{fides, qua creditur}) is indifferent to the object of faith (\textit{fides, quæ creditur}), that it, on the contrary, as shown above, is interested in it with infinite passion. This passion, that lies in the dialectic of certainty and uncertainty, is altogether personal. The believer must have, absolutely, certainty of the essential truth of the object of faith (--- the Spirit's efficacy in baptism; --- Christ's presence in the Supper); this certainty cannot be attained by the objective way: the assurance is in faith alone. Faith's assurance is subjective; and yet this subjective assurance is to take the place of objective certainty, and secure to itself objectivity: here is the tension. It is this tension, that sustains the passion and nourishes the inwardness: without an object of faith, no tension, and without tension, no true inwardness in faith. The sacraments are objective; but precisely for that reason are they, as objects of faith and objective
% --- p. 102 ---
acts, destined to be appropriated subjectively in the inwardness of faith. What is said here must not be misunderstood, as though it were now faith, subjective faith, that, in its inwardness, only fed upon itself, and nourished itself from itself alone; no, on the contrary! it is, indeed, precisely the object of faith, upon which faith feeds, it is the grace given in the sacrament, from which it draws its nourishment.

Should anyone now maintain, that the sacraments might, then, well fall away, as outward, objective acts, since faith itself, upon which everything depends, is inward and subjective: he thereby betrays, at the same time, that he has altogether misunderstood the nature and essence of inwardness, that is to say, of believing subjectivity. For the same reason, one then likewise misunderstands the nature and essence of the object of faith, insofar, namely, as one presupposes, that it, in order to establish its objective validity, must be treated and appropriated in objective knowledge. The matter stands the other way about: for objective knowledge, the object of faith vanishes altogether, that is to say: it transforms itself, and becomes an object of knowledge. As an object of knowledge, it then lies, in turn, under this contradiction, that it can be appropriated neither in faith nor in knowledge, but is only volatilized into an indeterminable something-between, a fantastic something, of which one does not rightly know what to say.

But in a dogmatics, that, although it does, indeed, say a good deal about subjectivity and inwardness, nonetheless does not let the problem of inwardness come forward at all, it is
% --- p. 103 ---

vain to seek any closer information concerning the relation of faith to its object.

As regards, now, in the first place, \emph{the relation between faith and historical objectivity}, there is, indeed, to be sure, a good deal said (Introd. §§ 22--23; 27--28, and elsewhere) about the material and formal principle of Protestantism, about holy Scripture and the Symbols, about sacred history and the witnessing Church; but if anyone, in this, for the rest, well-written presentation, expects to find a single decisive determination, he will, without doubt, be remarkably disappointed; for the fundamental question is lacking, approximation reigns, and the whole flows out --- into a fluid approximation. --- ``The objective canon for all Christianity,'' --- it says (§ 22) --- ``is, now, indeed, assuredly nothing other than Christ himself, in his holy, redeeming personality.'' --- Where, now, is Christ to be found? The nearest answer is the same as the Catholic one: in the Church! The right relation to the exalted Christ, glorified in the Church, is, again, conditioned by the right relation to the historical Christ, and thus we are led back, then, to the apostolic Church. ``For this stands firm: either it can no longer be determined, what Christianity is, in which case Christianity is no divine revelation, but only myth or philosopheme, \emph{or} there must be given a reliable tradition of the apostolic conception and appropriation of Christ, by means of which every succeeding period is able to preserve its connec-
% --- p. 104 ---
tion with the apostolic Church and genuine Christianity.'' --- Correct, altogether and entirely correct! but now the difficulty comes. On the view here set forth, Protestantism must, indeed, coincide altogether with Catholicism; so, then, the point is the deviation. The deviation from Catholicism is, namely, this, ``that we, with the Reformers, find the perfect, reliable form of the apostolic tradition only in the holy Scriptures of the New Testament, and that, therefore, Scripture is the last critical touchstone (\textit{lapis lydius}), that decides the Christianity of tradition.'' But is this objective \textit{lapis lydius} now also really altogether reliable, and in every respect sufficient? --- Not quite so! ``There is, indeed (§ 23), a conception of the scriptural principle, according to which nothing can have validity in the Church, except what can literally warrant its biblical origin. But this view is altogether foreign to the Lutheran Reformation, though traces of it are, indeed, to be found in the Swiss.'' --- But, for everything, that a Christian needs to know unto salvation, is, then, surely, the holy Scripture of the New Testament, the one right, objective, and critical touchstone (\textit{lapis lydius})? Yes, in one way, but yet not altogether so; ``for the regard for salvation is manifestly an individual regard, and the misapprehension of a truth, that, in one individual, may be no hindrance to salvation, may, in another individual, who stands on a higher level of the development of consciousness, become dangerous to salvation.'' --- But, as
% --- p. 105 ---
the fixed, infallible foundation for the dogmatic type of doctrine, holy Scripture must, surely, nonetheless, be able to lay claim to complete, absolute, and inalienable authority, and thus, in this sense, be, after all, the \textit{lapis lydius} both of dogmatics and of the dogmatician? --- Yes, in a certain way, but yet not altogether so. ``To be sure, Christian subjectivity (§ 24) is born out of the maternal womb of the Church, and must, from the first, find itself in an outward relation of dependence to the Church and to Scripture; but just as we, above (cf. § 8), showed, in general, that man's relation to God must be transposed, from a relation of dependence, into a relative relation of freedom: so, too, does this hold, in particular, of the relation to Christian revelation. Christian subjectivity must, in the course of its development, come to the point, where it no longer finds itself in the mere relation of dependence to the positively given, but in a free relation of exchange, a free relation of mutuality to it.'' --- And what, then, is further required, to bring Christian subjectivity to the point, where it no longer finds itself in the mere relation of dependence to the positively given? ``The outward canon'' (§ 24) ``points to an inward canon, as the principle for its understanding, that is, to the Christian \emph{regenerate consciousness}, in which God's Spirit bears witness with man's spirit (\textit{testimonium spiritus sancti})!'' --- See, here we are, again, at the point, the true \textit{punctum saliens}! but we must then ask, where this Christian regenerate consciousness
% --- p. 106 ---
is to be found? --- ``In the witnessing Church!'' --- This answer would be altogether satisfying, if it were only sufficiently shown, how the witnessing Church relates itself to the ecclesiastical witnesses. --- ``A witness can testify, that he has believed it (i.e. the historical fact, that God has become man), and thus, that it, far from being a historical certainty, runs counter to his understanding, but such a witness does, indeed, repel just as much in the same sense as the absurd, and a witness who does not thus repel is \textit{eo ipso} a deceiver, or a man, who speaks of something quite other, and such a witness can, then, help one to nothing, either, except to get certainty about something quite other. A hundred thousand individual witnesses, who, precisely by the peculiar character of their testimony (that they have believed the absurd), become individual witnesses, do not, either, \textit{en masse}, become something else, so that the absurd would become less absurd --- why? because 100,000 people, each for himself, had believed, that it was absurd? On the contrary, those 100,000 witnesses repel, again, altogether just like the absurd.'' Joh. Cl. --- ``Each one of us single individuals possesses faith only in a certain limited measure, and must, surely, guard against making our own individual, perhaps somewhat one-sided, perhaps even somewhat sickly, life of faith the rule for all believers. Only that great single individual, who, as the Apostle teaches us, shall, through the ages, grow up unto a perfect man, can fully realize the concept of the believer in the whole of its spirituality, the whole of its many-
% --- p. 107 ---
sidedness and rich multiplicity, can alone possess the fullness of faith and of all the gifts of grace.'' --- But if, now, ``each one of us single individuals'' possesses faith only in a certain measure, it follows, too, indeed, from this, that ``each one of us single individuals'' has Christian subjectivity, and the \textit{testimonium spiritus sancti}, only in a certain measure. Since, now, the outward, objective canon means nothing without the inward, and the inward is, in ``each one of us single individuals,'' only imperfect and relative: how, then, does it become possible for the Protestant Church to come to consciousness of itself, to secure to itself the testimony of the Spirit, the full and valid \textit{testimonium spiritus sancti}?

In order to limit the question, and direct the inquiry to a single, definite point: \emph{how does the evangelical-Lutheran Church get its adequate type of doctrine, its dogmatics?} The Church itself, considered as ``that great single individual,'' does not, nonetheless, after all, write the dogmatics. Dogmatics is written by a dogmatician. But every dogmatician is, indeed, a single individual; and yet this single individual is to represent, as dogmatician, ``the Christian regenerate consciousness, in which God's Spirit bears witness with man's spirit (\textit{testimonium spiritus sancti})'', is to present, though a single individual, the rule of faith for all believers, is, in other words, to identify his Christian subjectivity with the Church's true objective essence: how is this possible? could one not get to know, how the dogmatician manages to solve this task, as difficult as it is important? To abstract from one's singularity, and lose oneself in a historical-objective con-
% --- p. 108 ---
sideration of Scripture, tradition, and symbol, does not help, for, even when the objective presentation succeeds at its very best, one thereby, nonetheless, gets no further than to an outward conception of the outward canon, a factual determination of what Christianity is in its facticity; but such an objective presentation, that can, moreover, just as well be given by a non-believer as by a believer, still proves nothing about the objective truth of the Church's doctrine. --- ``The outward canon points to an inward canon, as the principle for its understanding''; dogmatics is (§ 2) ``not a merely historical presentation of what has been, or is, truth for others, without being so for the presenter himself''; but if this is so, we must, then, necessarily assume one of two things: either there are, in the Church, certain single individuals, who, in ``Christian subjectivity,'' possess the content of faith in so ample a measure, that they thereby become able to appropriate the whole doctrine of the Church, and expound the rule of faith for all believers; or else the evangelical Church must, in an infinite approximation, continue to seek its inward canon, without, however, ever attaining it in this approach, that is to say: without ever becoming conscious of itself in its ideal subjectivity. That, in the latter case, the Church can obtain no reliable dogmatics is self-evident; for when the dogmatician, as a single individual, acknowledges subjectivity only to a certain degree, and thus, in order to remedy this deficiency, must abstract from ``the single individual,'' and become objective to a certain degree: then the canon set up
% --- p. 109 ---
holds, too, indeed, only to a certain degree, and the evangelical-Lutheran dogmatics continues to go from approximation to approximation, without ever finding its promised πᾷ στῶ. If, on the other hand, the single individual undertakes to be a dogmatician, and, as dogmatician, that is, as ``the organ of his Church,'' at the same time to be ``the organ of his science'': then he must, too, scientifically or unscientifically, give due account of ``Christian subjectivity.'' By being a little subjective, as the single individual, and, for the rest, objective, as the organ of his Church, he accomplishes nothing. Only by himself centering in Christianity's absolute principle of subjectivity, and showing, how the single individual, not as this or that accidental individual, but as the living organ of the idea and the spirit, is able to appropriate the factual revelation, can the dogmatician develop the truth of Christianity, and vindicate its objectivity. If, on the other hand, one will be subjective only to a certain degree, only call subjectivity to mind in certain subjective assurances, in order, all the more freely, to be able to fly over ``the single individuals'' by a subjectless hovering in the historical-ecclesiastical objectivity, one misses the problem of faith. But if the problem of faith is missed in the doctrine of faith, what use is it, then, to want to stand on πᾷ στῶ, and be objectively distinguished, and ecclesiastically high and mighty, and behold --- the glory of God in speculation?

The older Lutheran dogmaticians used, for the many-sided conception of faith, to bring out a distinction, that it greatly astonishes me, Dr. M. has not touched upon at all, namely
% --- p. 110 ---
the distinction between the human and the divine side of faith (\textit{fides humana} and \textit{fides divina}), a distinction, that, with some modification, corresponds roughly to the distinction between historical and religious faith (\textit{fides historica} and \textit{fides religiosa}). By occupying oneself with dogmatics' relation to holy Scripture, one comes, involuntarily, to reflect on this distinction, and thereby, again, to consider dogmatics' relation both to introductory criticism and to modern hermeneutics. To be sure, a dogmatician ought neither to enter into the special inquiries of criticism, nor the particular rules of hermeneutics; but insofar as it is a matter of securing faith, he ought, surely, nonetheless, to see himself able to draw a boundary line, and set up a canon, both for introductory criticism and for hermeneutics. --- Why the research of criticism, insofar as it keeps to a historical conception of historical data, cannot be rejected in more recent Protestantism, needs here no closer explanation. But, as the critical introduction now continues its researches, and these researches, in so many different ways, change and modify the outward canon, and thereby, too, change and modify the historical conception (\textit{fides historica}): how, then, does it come about, that the dogmatician can still use the holy Scripture of the New Testament as an inspired Scripture, how does the historical faith relate itself to the religious faith? But on this point, ``the Christian Dogmatics'' contains no definite information; for what (p.
% --- p. 111 ---

56--58) is said about the one-sided application of the formal principle in the orthodoxy of the seventeenth century, and about the natural reason of rationalism in the eighteenth century, decides nothing as yet. One can, indeed, readily concede to the author, that rationalism was nothing other than ``a great \textit{theologia irregenitorum}, that flooded Protestant Christendom,'' if he will, then, in return, instruct us, how, then, one is to secure oneself a \textit{theologia regenitorum}. --- That hermeneutics, as a philological science, may, in many respects, be entitled to distinguish between spirit and letter, content and form, no one, surely, will deny, who has any conception of what rational interpretation means; but, when the talk is, now, of a holy Scripture, that is set up as a \textit{lapis lydius}, an inspired Scripture, that not merely contains God's word, but also is God's word: ought it not, then, to be the dogmatician's task to indicate a \textit{criterion}, by which one could, with tolerable certainty, judge, when and how far that distinction of spirit and letter is for or against faith, and thus more closely determine the relation between \textit{fides humana} and \textit{fides divina}? Since, however, I am, here, neither writing a dogmatics, nor an introduction to a dogmatics, I cannot, in this context, enter into any special inquiry concerning it, but will only remark, that ``the Christian Dogmatics,'' far from affording any guidance, does not even seem to suspect the problem.

``The concept of sacred history'' --- it says
% --- p. 112 ---
(p. 24, § 16) --- ``is inseparable from the concept of the miracle;'' after having made some cosmic-preliminary observations on the miracle of the incarnation and of inspiration, the author then proceeds, in §§ 17 and 18, to treat of ``the opposition between supranaturalism, on the one side, and naturalism and rationalism, on the other.'' What is said, in these introductory paragraphs, about the nature of natural law and the teleological character of the miracle, is, in reality, roughly the same as what one can otherwise find in a simple ontological compendium\footnote{Cf., e.g., \textit{The Propaedeutic Logic}, pp. 144--157.}; the difference is only this, that ontology seeks to lay out its development upon a little dialectic, while ``the Dogmatics,'' on the other hand, can express itself quite thetically; but this is, indeed, only a difference of form.\footnote{As an example of such a difference of form, one may compare what is said about the personality of the angels in ``the Christian Dogmatics'' (§ 69), with what is said about personifications and intermediate hypostases in \textit{The Propaedeutic Logic} (pp. 182--86).} The actual development of ``the concept of the miracle'' the author has reserved for ``the dogmatic system itself''; but it would, nonetheless, have been highly desirable, that he had not done this. For, before one can think of a systematic development of ``the concept of the miracle,'' one must, surely, first, in the introduction, have secured \emph{faith's relation to metaphysical objectivity}, and assured oneself, how far the miracle, now, can really let itself be developed in ``the system,'' and mediated in ``the concept.'' Had Dr. M., in full earnest,
% --- p. 113 ---
put to himself the question: how does the natural relate itself to the supernatural in Christian revelation, and \textit{in specie} in sacred history?\footnote{On the significance of the miracle in sacred history, cf. \textit{The Gospel Faith and the Modern Consciousness}, pp. 17--20; 87--92; 351--55; 401--405, and elsewhere.} then I am convinced, that he, according to his own confession, not for any price wishing to be at variance with the believer, would at once have let speculation go, and given us a quite different development of ``the dogmatic system itself.'' --- That ``sacred history'' cannot be conceived directly, or by analogy with any straightforwardly historical fact whatsoever, an admissible, though indirect, proof has, so far as I can judge, been furnished in the Pseudonym's \textit{Philosophical Fragments}.\footnote{Published 1844.} The opposition between sacred and profane history is not yet expressed by speaking of the former as the most exalted and divine thing conceivable; not by calling sacred history ``the sinless,'' and world history ``the sinful history''; for the question itself turns precisely on the \textit{conditio sine qua non} for the conception of this opposition. The historical Christ is, and remains, as long as this world stands, an offensive contradiction\footnote{Christ is a sign that is spoken against. \textit{The Gospel Faith and the Modern Consciousness}, pp. 288--289.} for the direct historical view,
% --- p. 114 ---
and every attempt, by the way of objective knowledge, whether historical or metaphysical, to remove this contradiction, is only misunderstanding, error, and folly. Least of all, then, can one remove the contradiction, when one does not even let it come so far into view, that the problem shows itself.

So long as the determination of faith's relation to objective knowledge is kept in ambiguity, the dogmatic cognition can by no means assume any fixed, definite character. It must, surely, in truth, be capable of being decided, how it really stands with the content of faith, whether this content is something, that is to be comprehended, or not. Let us choose the dogma of the Trinity: is this dogma, then, determined to be a philosopheme, and expounded in objective comprehension? In order to avoid ``the fluid,'' one had best give a definite answer: Yes or No. If one answers Yes! then faith, that is to say, the simple faith, is only a provisional makeshift, that one makes use of, until one can attain that, which is both higher and surer and more glorious, namely the objective concept. Even if one cannot at once attain the perfect concept, one must, nonetheless, for the time being, make do with the less perfect, for faith, indeed, will comprehend, the dogma, indeed, will absolutely be transformed into a philosopheme. But what, then, follows from this? --- ``That thou shalt, in the first place, believe in the triune God, though thou dost not yet fully comprehend his essence, though thy cognition is, as yet, only a shadow-cognition (cf. Dogm. § 55).'' --- But this answer is
% --- p. 115 ---
ambiguous. When I am told, that I am to believe in the triune God, though I do not yet fully comprehend his essence, it is, nonetheless, presupposed, that I have a shadow-comprehension, and that this conceptual shadow shall, then, for the time being, be the foundation, on which my faith is to rest. But suppose, now, that I, with my comprehending, had come to the cognition, that the ecclesiastical conception of God's Trinity conflicted with the concept: what should I, then, do? --- ``Yes, then thou shouldst, nonetheless, continue to believe in the triune God; for the doctrine of God's Trinity is, when one rightly considers, nonetheless, no proper philosopheme after all, but a dogma.'' --- Good, very good! But tell us, now, in which of these three cases, then, is my faith best secured and grounded, either: when I believe \emph{``with the Concept,''} or: \emph{without ``the Concept,''} or: \emph{against ``the Concept''}? --- ``Naturally, in the first case.'' --- Should it, then, really stand so? Let us weigh the matter, and, as the Pseudonym says, assume first the one, and then the other. Assumed, then, that faith is to be helped by knowledge, and, so to say, cast a part of its burden upon ``the Concept,'' then it must, indeed, be far easier to believe with ``the Concept'' than without ``the Concept''; the higher ``the Concept'' rises, the easier faith becomes, the more its exertion diminishes. It goes, in this case, with faith as with the poor laborer, who, by the works of his hands, earned his daily bread, until fortune once
% --- p. 116 ---
came to look upon him, and visit him with an inheritance --- ``from the rich uncle''! So long as the man of no means, left to himself, labored early and late, his strength grew with his labor, and he thanked God; for there was, indeed, otherwise no one in the world, who would labor for him. But, when he now received the inheritance, and came into prosperity, and ceased to labor, since he ``no longer needed it,'' but could ``keep people for himself,'' and let others labor: what then? Well, then he began, precisely, to feel himself weak; and his courage sank, the more his wealth grew, and his strength diminished, the more familiar he became with the ``good days,'' and with the company that most readily gathers in the ``good days,'' that is, with ease, idleness, and tedium. Roughly so, and, for the most part, not otherwise, does it seem, precisely, to stand with faith, the faith, that, by a stroke of ingenious good fortune, becomes heir to knowledge. So long as the simple faith, left to its own exertion, watchful, early and late, must labor without ``the Concept,'' its strength grows with its labor, and it thanks God; for there is, indeed, otherwise no one in the world, who will labor for it. But, when, then, philosophy, not unlike that rich uncle, once, as on its deathbed, bequeaths it a rather considerable capital of ontological ``shadow-cognitions,'' and thereby puts it in a position to keep concepts for itself, and let the objective thoughts take the brunt: what then? Well, then faith will, unfortunately, only too easily be weaned from exertion, weakened in idleness, and sicken in softness.
% --- p. 117 ---
Let us now assume the opposite, namely, that faith, in beginning without ``the Concept,'' learns to rely upon itself in the face of objective knowledge: would it, then, yield a hair's breadth, even if objective knowledge should once get the dogma transformed into a philosopheme, and the philosopheme, in turn, reduced \textit{in absurdum}? By no means! In such a conflict, faith will, precisely, become fully conscious of itself as faith; in such a struggle, faith, in trust in Him who called it forth out of nothing, in humble trust in the Spirit that bade it be, in this world, what it is, faith and nothing other than faith, will gather its strength, and lift up its head, and bear its burden. The proposition: I believe, \emph{although} I cannot comprehend, is ambiguous, and can first be rightly construed in the proposition: I believe, \emph{because} I cannot comprehend; \textit{credo, quia absurdum est}.

That the old Tertullianic proposition, \textit{credo, quia absurdum est}, looks, at its first appearance, somewhat harsh and barbarous, and, therefore, can scarcely be welcome to speculation, shall never be denied. Meanwhile, one is, nonetheless, greatly mistaken, if one regards it, without further ado, as an arbitrary scandal, a raw stumbling-block, that a barbarian has cast into the Church, in order to mock science and scoff at coherent thinking. To be sure, it is a stumbling-block, but this stumbling-block is, at the same time, to be regarded as an incorruptible sentinel, a strict and faithful border-guard, who stands and watches, where faith and knowledge hold their meeting, that they may not, in mutual infatuation, outwit one another, and cheat one
% --- p. 118 ---
another out of something. Whenever, now, either knowledge trespasses upon faith, and is on the point of giving up what is its own, in order to appropriate to itself what belongs to faith, or faith, enamored of knowledge, forgets its dignity for an imagined good fortune with speculation, the old border-guard stands by, and addresses them sternly, and parts them with the harsh word: \textit{credo, quia absurdum est}.

Only in relation to the absurd can it rightly become manifest, what faith is, what it means to believe. For speculative philosophy of religion, God's essence is, and must be, comprehensible; for here the point is to bring out and hold fast the conceptual side of the divine. Speculation's God I am to comprehend; but precisely because I am to comprehend, I am not to believe. To conceive speculation's God as the God of faith is a naive misunderstanding, just as it is a naive misunderstanding to wish to believe --- in speculation. For dogmatics, on the other hand, God's essence is, and must be, incomprehensible, since it is to be the object of faith and worship. When the dogmatician, nonetheless, accentuates the conceptual side, and conceals the incomprehensible in philosophical turns of phrase and speculative circumlocutions, one may always be sure, that he has failed his dogmatic task, and makes his point falsely. --- This must not, however, be understood, as though the dogmatic presentation were to be devoid of concepts, and dogmatics afraid of any possible contact with the philosophy of religion. The proposition, ``God is incomprehensible,'' can, namely, be construed in a highly different way, according as it takes its
% --- p. 119 ---

direction toward knowledge, or toward faith. In relation to knowledge, no one can appeal to God's incomprehensibility, without his also having to render science an account, of how far he has, indeed, also cognized the comprehensible. But those, in number, who, in school theology, most readily take up the word for God's incomprehensibility, not because they really wish to uphold faith, but because they thereby spare themselves the trouble of a tolerably thorough metaphysical preliminary study, are, naturally, always the true popular majority; and speculation was, assuredly, well within its rights, when it, under the firm of \textit{vulgus odi}, took upon itself to chastise vulgar rationalism. But, in relation to faith, the proposition, that God is incomprehensible, properly speaking, needs no scientific proof at all, since it can, indeed, straightforwardly be said to be grounded in the very essence of faith. The believer need not, precisely, be initiated into the mysteries of speculation, in order to be able to meet any philosopher, who claims to have comprehended the God of revelation, with the objection, that one, as regards this point, must, assuredly, be in an error. But, insofar as a dispute can, now, arise about this, between the believer and the philosopher, it is of importance, that they mutually understand one another, which then, again, requires, not merely that the philosopher acquaint himself with faith, but that the believer, too, know how he is to operate with the concept. A provisional agreement, on the concession, that one can, thus, half and half comprehend, and, therefore, too, only half and half need to believe, is a
% --- p. 120 ---
middling result, that satisfies only ``mediocrity.'' If the inquiry is to be carried through, the philosopher must continue to drive the demand of the concept to the utmost, and the believer, who follows along, and understands the philosopher to the last dot, must, every time, protest the demand with the turn of phrase: ``I comprehend, that I cannot comprehend.''

\emph{I comprehend, that I cannot comprehend}: on this basic proposition, dogmatics can freely meet with the philosophy of religion, and appropriate its yield, without paying for it with its independence. It is said, in the Gospel, of the sinful woman, who fell down before the Savior, that she broke her alabaster jar, and poured out the precious ointment over his feet. She loved much, and therefore she also wished, right gladly, to sacrifice much, to sacrifice the most costly thing she had. But, in order to sacrifice the most costly thing, she must, indeed, surely, first have acquired the most costly thing, must first have bought both the ointment and the jar. Let this woman, who, in relation to existence, is, indeed, always and remains, for the believer, a model,\footnote{Cf., e.g., Dr. Martensen's sermon, ``Mary Magdalene,'' \textit{Sermons}, 1847.} now, too, in relation to knowledge, be a model for the dogmatician. Let us but concede, that knowledge and science are precious goods, yes, let us even, as philosophizing beings, regard speculation as the most costly thing; but as dogmaticians, we must, nonetheless, not forget
% --- p. 121 ---
ourselves, must not forget, that even if knowledge were, of all human things, the most costly, that it is precisely then the most costly thing, that is to be sacrificed.

In order to illustrate this point with a definite example, I will here merely direct attention to a section of Dr. Martensen's Dogmatics, that, in my opinion, is among the most successful and best carried through, I mean the section, that treats of ``the Lord's Last Coming, and the Consummation of All Things.'' This section, in which the testing thought moves with far greater freedom and elasticity, than is the case in most of the preceding sections, ends, too, quite rightly, with an antinomy, that cannot be resolved, and cannot be mediated in any concept, namely the antinomy between universal apokatastasis and eternal damnation. It makes a strange, almost touching impression, to meet this strict, terrible antinomy here, right at the close of the work. It is, indeed, almost like a symbol, that the dogmatician's path through the system shows itself as an image of a seeking human being's path through life; for, in the life of a human being, who seeks after the truth, it happens, at times, that the seeker is, here and there, quite near to finding, what he seeks, but, nonetheless, constantly confused by distracting side-glances, and misses what is sought, until, at last, in a final hour, when the long way has been traversed, he has his eyes opened, and sees the watchword, with which he ought properly to have begun. Roughly so, it seems, has it gone with the author of ``the Christian
% --- p. 122 ---
Dogmatics.'' He begins with the intention of wishing, as far as possible, to mediate ``the works of the Most High'' (§ 33); but see, after he has, then, throughout the whole book, here and there, attempted to comprehend and mediate, as far as possible, he ends with the result, that the very last antinomy, the antinomy, in which the whole inquiry, at last, concentrates itself as in its knot-point, does not let itself be mediated, but strikes down every concept, and stands fast (§ 288) as --- ``a cross for thought.'' But if we must now, indeed, regret, that this, the dogmatician's last word, has not also become his first word, we nonetheless rejoice at the thought, that he does not yet stand at the end, either of his theoretical or his practical career, and console ourselves, that what has not already happened may, indeed, in time, happen, since the man is still in his strength, and lacks neither talent nor ability, if he could only find the problem.

But that Dr. M. has not yet found the problem, not even in the last section, is only too clear from his ``coherent cognition.'' The antinomy remaining at the close of the system arises, namely, not so much from a dialectical insight into the essence of the dogma, as from an outward glance at the positive doctrine of the Church. ``Is there,'' --- it says (§ 183) --- ``an eternal damnation, or may one assume a final conversion of the damned, a universal apokatastasis, a restitution of all fallen beings, such that God, in the full sense, becomes all in all? The Church has never wished to enter into this latter''; --- \textit{ergo} the dogmatician, who
% --- p. 123 ---
is only the organ of his ``science,'' in that he is, at the same time, the organ of his ``Church,'' will not, either, enter into it. Had the dialectical moment, that hides itself in this antinomy, been discovered and held fast, the dogmatician would at once have realized, that it has, Christianly understood, no greater difficulty with an eternal damnation, than with an eternal blessedness,\footnote{When two thoughts stand in an indissoluble relation to one another, so that, if one can think the one, one can, \textit{eo ipso}, think the other, it happens, not infrequently, that there goes, from mouth to mouth, from generation to generation, an opinion, that makes it easy to think the one thought, while an opposite opinion makes it difficult to think the other, indeed, establishes the practice, of being skeptical in relation to this one. And yet the true dialectical relation is, that he, who can think the one, can, \textit{eo ipso}, think the other --- if he has thought the one. I refer, hereby, to the quasi-dogma of the eternity of the punishment of hell. The problem, the \textit{Fragments} set forth, was, how something historical can become decisive for an eternal blessedness. When one says ``decisive,'' one says, \textit{eo ipso}, that, in that blessedness is decided, unblessedness, too, is decided, be it now as included or as excluded. The former is now supposed to be easy to understand, every systematician has thought it, and we are, indeed, all believers, it is child's play to get a historical point of departure for one's eternal blessedness, and get it thought. In the midst of this security and reliability, there occasionally comes forward the question of an eternal unblessedness, decided by a historical point of departure in time: see, that is difficult, one cannot come to agreement with oneself, about which one is to assume, and one comes to agreement with oneself, to let it stand as something, that one
% --- footnote continues on p. 124 ---
can, on occasion, use in popular lecturing, but that is undecided --- that, that, that, and so it is, indeed, decided; nothing is easier --- if one has decided the first. Wondrous human thoughtfulness, who can gaze into your thoughtful eye without a quiet uplifting!'' Johannes Climacus.} and that, accordingly, the incomprehensibility of this one point is sufficient to prove the incomprehensibility of the whole of Christianity.

The Christian dogmas all stand in the same relation
% --- p. 124 ---
to objective knowledge; the one is exactly as incomprehensible as the other. If one is on the way to comprehending a single dogma, then one is on the way to comprehending them all; but if there is a single one, that contradicts the concept, and stands fast as --- ``a cross for thought,'' then this is an unmistakable sign, that the dogmatician must, too, with every one of the rest, strike a cross for thought. That an unsystematic mind should not at once realize this is not to be wondered at; but that Dr. M., who is, after all, indeed, a systematic mind, which can, among other things, be proved by this, that he looks down so loftily upon the single individuals, who can only think in ``scattered thoughts and aphorisms, whims and flashes,'' that Dr. M. has not realized this, is very strange. Take whichever systematician you will: it shows itself at once, that he, as regards the cognition of the divine, either aims at a comprehending, or a not-comprehending, so that his last word always corresponds to his first. Kant took upon himself to set up the antinomies, but, since he could not resolve them all, he resolved none at all; for Kant was a systematic mind. Hegel combated and overcame Kant, or at least meant to overcome him, in that he
% --- p. 125 ---
resolved, not a single antinomy, but all antinomies; for Hegel was a systematic mind. But, it goes without saying, every systematic mind, whether he is, for the rest, philosopher or theologian, must, indeed, too, from the very first, be certain of his problem; for, so long as one has not yet found the problem, it, nonetheless, all the same, has no bearing whatsoever on systematic thinking, and ``coherent cognition.''

\begin{center}
\emph{d) ``The Christian Dogmatics'' is without principle.}
\end{center}

It would, assuredly, become a very extensive labor, if an exacting critic were to undertake to review ``the Christian Dogmatics,'' point by point, in all its joints and connections; perhaps the one difficulty would then continue to develop out of the other, so that the criticism would, in the end, find no end at all. That would, perhaps, well be possible. Meanwhile, this examining review cannot, nonetheless, lay any particular weight upon a great number of individual objections, since it presupposes, far rather, that the work before us is, in every respect, so well thought through, and so defensibly worked out, that there is, at bottom, nothing whatsoever to be objected against it, except this one thing, that the system is without principle. If one will but overlook this one point, then all the rest is, assuredly, wholly and entirely, as it ought to be; one will, then, in the work, considered as a whole, have ample occasion to acknowledge both ``the definite'' and ``the clear,'' as well as ``the
% --- p. 126 ---
spirited'' and ``the profound,'' that, in so many passages, and under so many turns of phrase, speaks to him, who has any ``deeper sense''; I, for my part, should then, at least, not have concerned myself with any objection, but conducted myself just as handsomely acknowledging as most others. But since I, as already noted above, have this time resolved to look at nothing else, than solely the problem and the principle, I should belie my innermost conviction, if I did not now, publicly and openly, avow the claim, that ``the Christian Dogmatics'' is without principle. Should anyone wish to object to this, that this claim is altogether too invidious, since ``the Dogmatics,'' far from suffering from any total lack of principle, might, rather, be said to have a certain superabundance of different principles: I shall have nothing against this, if by a principle one understands merely the same thing as a basic proposition. For no one will, surely, deny, that there are to be found, in this Dogmatics, a multitude of basic propositions, concerning ``speculation'' and the religious ``relation of existence,'' that are, in principle and essence, so dissimilar, that a highly questionable revolution would, assuredly, arise in ``coherent cognition,'' if these heterogeneous basic propositions really came to word, and were allowed to break through. In order to avert this danger, and forestall so sorry a mishap, Dr. M. has already, in the introduction, been intent upon disciplining these propositions, and taking such disciplinary-restrictive measures on both sides, that
% --- p. 127 ---

the one basic proposition should not offend the other. By the application of a discipline as gentle as it is well-calculated, the author has succeeded in bringing about a certain balance, calm, and order, in the whole work; nonetheless, this apparent calm is not, after all, wholly to be relied upon; for the principles smolder, and the hostile propositions cast challenging side-glances at one another. --- When, now, ``faith'' once gathers all its own, and ``speculation'' all its own: how, then, shall these two ideal powers come to an understanding with one another, and reach a settlement? Even if one wished to mediate between them, and, for the Church's sake, grant ``faith'' twice as great an influence as ``speculation,'' the system would, indeed, nonetheless, at the very least, have a principle and a half; but a dogmatics on a principle and a half is without principle.\footnote{One knows, how many different, mutually contradictory schools formed themselves after Socrates, because he was not a schoolman \textit{ex officio}; but one does not yet know, how many highly different, mutually contradictory dogmaticians will, in time, form themselves out of Dr. M.'s Dogmatics, perhaps for the opposite reason, that this Dogmatics is a chief dogmatics \textit{ex officio}.}

See, that is why it is, that I, this time, hold to the question of principle, and will, so far as I am able, continue to do so, even if the man, whom I esteem as few others, and would most gladly see as my friend, should, in this dispute, appear as my opponent. I have, at least, accepted his challenge, though it did not, literally,
% --- p. 128 ---
run upon my name, and shall now, in closing, only add a couple more words about the character of the dispute, in order thereby, at the same time, to settle the provisional \textit{status quo}.

To keep a dogmatic dispute over principle so pure and unmixed, that nothing creeps in, that could, in any way, be construed as a personal offense to the other party, is very difficult, if not altogether impossible. Meanwhile, something is, nonetheless, gained, when the disputant, who has confined himself solely to the matter, guards against having those allusions, that, in the heat of polemic, may happen to strike the mistaken standpoint, and the illusions connected with it, construed, without further ado, as offensive attacks upon the opponent's person.

The dogmatic dispute over principle, that is here carried on, is by no means any immediate and inspired contest of talent. For such a contest, one who philosophizes, who must continually work over his own and others' thoughts (without, however, thereby meaning to claim rank with genuine genius), can, naturally, feel neither desire nor courage. In an immediate contest of talent, it may be allowed, that, when the one is best with the subjective, the other best with the objective, each does best to follow his own genius. But, when the talk is of a principle, and namely of a religious-ethical principle of cognition, the matter stands otherwise. --- From my very first acquaintance with Dr. Martensen, I have retained the impression, that this man was excellently gifted to become a dogmatician, and I acknowledge, still, that our university could not
% --- p. 129 ---
do without such a dogmatician, without too keenly felt a loss to theology. But it is, nonetheless, my conviction, that Dr. M., relying on his dogmatic talent, has not been sufficiently open to hints and objections, that might, otherwise, perhaps, have warned him against several missteps in objective knowledge; that he has, with all too great a certainty, remained standing on the standpoint once given him, and has, precisely thereby, come to overlook the problem, and misjudge speculation. But, when it is now willingly conceded, that it is neither talent nor powers, that our dogmatician lacks: should one not, then, without giving offense, be able to dispute with him, about how these powers ought, preferably, to be applied?

Our dogmatic dispute over principle is, happily, no hateful dispute over heresy, either. If the single individual will not exalt himself to the ideal of faith, and Dr. M., for no price, will be at variance with the believer: then this point will, indeed, settle itself, as though of its own accord. What is, then, said in ``the Preface'' about the single individuals, who wish ``to make their own, perhaps somewhat one-sided, perhaps even somewhat sickly, life of faith the rule for all believers,'' we will now take in the best sense, and write it clean off in the book of oblivion, that we, instead of the passion of faith, that ethically transforms life, should not awaken quite another, and very unpleasant, passion, namely that which recalls \textit{rabies theologorum}.

So, then, is our dogmatic dispute over principle, too, no egoistic contest of rivalry, either, as though it were now a matter
% --- p. 130 ---
of who should occupy the first place, either in science or in the Church. As for myself, I know, of myself, that I am not organized to take the lead, and feel, therefore, not at all troubled by the ambition to wish to be the first. In the important matter, of which we treat, I have not, moreover, either, invented anything (if there is anyone, who has here seen and found something, it is, no doubt, someone else, and not I); but even if I cannot boast of any significant or ingenious productivity, I will, nonetheless, always, to the best of my ability, strive to keep pace, and shall be very loath to give up --- coherent cognition.

Should anyone suppose, that this sincere confession is, nonetheless, all the same, only a spuriously modest irony, that sufficiently betrays its origin, since it stands in so glaring a disproportion to the presumptuous tone, in which this review has, presumably, attacked Dr. M.: then I must, first, refer him to the explanation of the eighth commandment in Luther's Small Catechism, and add, besides, a word. A review, critique, or notice is not always just as fair, in proportion, as it is moderate in its objections. I have no intention whatsoever of taking any ``free license''; but I hate those half-baked, moderate reviews, that, all too envious to grant an author what rightly belongs to him, and all too cowardly to dare raise any decisive objection, only mystify opinion by their insipidity, and make it difficult, if not impossible, for the party concerned to reply.
% --- p. 131 ---
The more definitely criticism sets forth its objection, the easier it is for the author to conduct his defense; and it will, perhaps, be granted me, that I have, in this piece, striven to make it as convenient as possible for Dr. M.

Here, now, so far as I can judge, only one of two cases can occur: either I am, in the main, right, or else altogether wrong; but, in either of these cases, the dogmatician's position is, at the very least, just as advantageous as mine. If Dr. M. can refute my claims, he will, indeed, surely, win a victory, not merely over me (for that means little), but possibly, even, over every single individual, who, as ``the Preface'' has it, can only think in ``scattered thoughts and aphorisms, whims and flashes.'' If, on the other hand, he does not find it fitting to enter into any further dispute (since such a dispute might, possibly, give rise to quackeries), he need not, even then, be at a loss: he can, indeed, simply keep silent. With me, it is quite another matter. Should I even be found to be right, I nonetheless win no victory, since I do not, here, originally, offer anything of my own, but only appropriate, what I believe I have learned from another. If, on the other hand, I am found to be wrong, and am convinced of having misunderstood both the one and the other, I shall, then, indeed, have to bear the shame alone, since no one else, surely, can be served by taking over my misunderstandings, and making amends for my missteps. I have not, then, even the refuge of being able to keep silent. For, if it is true, that it is not he who strikes, but he who strikes back, who begins the dispute, then
% --- p. 132 ---
this, my self-defense, is, nonetheless, in a way, an attack, and I must, as the attacker, surely, know, what I owe my opponent. Should Dr. Martensen, then, prove to me, that I have really misunderstood his ``Christian Dogmatics,'' in its principle as in its problem: then I shall, indeed, assuredly, not attempt, either, to shelter myself behind any ambiguous silence, but shall, freely and openly, God willing, meet him with my answer, even should this answer consist only in the three words: ``I have erred.''

\bigskip
\noindent Copenhagen, 18 September 1849.
\hfill {\bfseries R. Nielsen.}\hspace{2em}

\begin{center}---\end{center}

\end{document}
